\documentclass{amsart}
\usepackage{amssymb,mathtools,mathrsfs,array,longtable}
\usepackage[hidelinks]{hyperref}
\hypersetup{
  pdftitle={Relations among exponential integrals, Bessel and Kummer functions},
  pdfauthor={Ruiyi Zhang},
  pdfsubject={Primary 11J91; Secondary 12H05, 33C10, 33C15},
  pdfkeywords={algebraic independence, E-functions, differential Galois theory, exponential integral, Bessel functions, confluent hypergeometric functions}
}
% Shared notation. Ordinary amsart dimensions and font sizes are retained.
\DeclareTextCommand{\textsection}{OT1}{\UseTextSymbol{OMS}{\textsection}}
\allowdisplaybreaks[1]
\makeatletter
\g@addto@macro\@adjustvertspacing{%
  \setlength{\abovedisplayskip}{1pt plus 1pt minus 1pt}%
  \setlength{\belowdisplayskip}{1pt plus 1pt minus 1pt}%
  \setlength{\jot}{1pt}}
\makeatother
\normalsize
\setlength{\emergencystretch}{2em}
\newtheorem{theorem}{Theorem}
\newtheorem{explicitestimate}{Theorem}
\renewcommand{\theexplicitestimate}{4}
\newtheorem{lemma}{Lemma}[section]
\newtheorem{proposition}[lemma]{Proposition}
\newtheorem{corollary}[lemma]{Corollary}
\theoremstyle{definition}
\newtheorem{definition}[lemma]{Definition}
\newtheorem{example}[lemma]{Example}
\numberwithin{equation}{section}
\numberwithin{table}{section}

\newcommand{\Q}{\mathbb{Q}}
\newcommand{\Qbar}{{\overline{\mathbb{Q}}}}
\newcommand{\Z}{\mathbb{Z}}
\newcommand{\C}{\mathbb{C}}
\DeclareMathOperator{\Ein}{Ein}
\DeclareMathOperator{\trdeg}{trdeg}
\DeclareMathOperator{\Gal}{Gal}
\DeclareMathOperator*{\Res}{Res}
\DeclareMathOperator{\GL}{GL}
\DeclareMathOperator{\SL}{SL}
\DeclareMathOperator{\PGL}{PGL}
\DeclareMathOperator{\Lie}{Lie}
\DeclareMathOperator{\rank}{rank}
\DeclareMathOperator{\Span}{span}

\makeatletter
\newcommand{\equationalias}[1]{\begingroup
  \protected@edef\@currentlabel{\p@equation\theequation}\label{#1}\endgroup}
\makeatother


\begin{document}

\title[Exponential integrals, Bessel and Kummer functions]{Relations among exponential integrals, Bessel and Kummer functions}
\author{Ruiyi Zhang}
\address{Independent Researcher, Columbia, Maryland, USA}
\email{ryzfsu@gmail.com}
\date{November 17, 2026}
\subjclass[2020]{Primary 11J91; Secondary 12H05, 33C10, 33C15}
\keywords{algebraic independence, E-functions, differential Galois theory, exponential integral, Bessel functions, confluent hypergeometric functions}

\begin{abstract}
We determine algebraic relations among finite families of special functions.
The families contain rational-exponential primitives and iterated entire exponential integrals.
They also contain Bessel functions of rational order and Kummer functions.
The Kummer functions have nonintegral first parameter and second parameter one.
We give the transcendence degree and generators of the relation ideal.
Explicit maps for normalized entire functions complement Gorelov's 2019--2021 classification of second-order equations.
We then specialize the relations at nonzero algebraic points.
Over explicit base fields, we obtain a lower bound for the transcendence degree of mixed
connection values together with Euler's constant, logarithmic constants and,
when required, pi.
We also prove a linear independence estimate with explicit constants for all sufficiently large heights.
It concerns one, the exponential of one, the reciprocal of that exponential and the entire exponential integral at one.
The qualitative results use differential Galois theory and E-function specialization.
The estimate uses integer polynomial approximations and derivative selection.
\end{abstract}

\maketitle

\section{Introduction}\label{sec:introduction}

The transcendence degree of Siegel's $E$-functions \cite{Siegel1929}
is preserved under specialization at nonzero algebraic ordinary points \cite{Shidlovskii1989}.
Beukers lifts relations \cite[Theorems~1.1--1.3]{Beukers2006}
using Andr\'e's $E$-operators \cite{Andre2000I,Andre2000II}.

Parameter coincidences, Wronskians and shared logarithmic companions
produce mixed relations. We determine their ideal, retaining normalized
additive constants. Nonclosed lists are treated by closure, specialization and elimination below.

\subsection{A first example}
Take $c\ne0$ and the linked blocks $\operatorname{Bess}(c,0)$ and
$\operatorname{Kum}(2c,1/2)$. They form one linkage class and leave three independent coordinates
over $T=\C(z)(e^{cz})$ and four over $\C(z)$.
Supplementary Section~\ref{sup:first-example} gives the maps, determinant, series and connections.

\subsection{The main results}\label{subsec:normalization}
The following data conventions define all terms in Theorem~\ref{thm:main}.
Work over $k=\C(z)$ with $D=d/dz$ and a finite admissible coordinate list
$\mathcal F$, including repeated input slots. Fix a simply connected
$\Omega_{\rm an}\subset\C^\times$ avoiding the finite poles of the data,
a point $z_*\in\Omega_{\rm an}$ and a branch $L=\log z$ there.
All functions are meromorphic on this domain, with differential constant
field $\C$. Set $\phi_{\mu,p}=e^{\mu(z-p)}/(z-p)$,
$\Phi'_{\mu,p}=\phi_{\mu,p}$, $\Phi_{\mu,p}(z_*)=0$ for
$(\mu,p)\ne(0,0)$, and $\Phi_{0,0}=L$.
A prime on a coordinate denotes $D$; on a function of its own argument
it denotes that derivative, so $D[f_a(sz)]=sf'_a(sz)$.

The coordinate dictionary in Section~\ref{subsec:blocks} specifies five types:
exponentials $e^{\lambda z}$; specified first-level primitives of finite sums
$\sum_\mu q_\mu e^{\mu z}$ with $q_\mu\in\C(z)$ and arbitrary complex
additive constants, together with the finite towers $\Ein_k(\eta z)$
($\eta\ne0$, $1\le k\le h_\eta$) of \eqref{eq:ein-tower};
Bessel blocks of nonzero scale $c$ and rational order, with the four
coordinates in \eqref{eq:bessel-zero}--\eqref{eq:bessel-nonresonant};
Kummer blocks $(F,F',\Sigma,\Sigma')$ of nonzero scale $s$ and
$a\notin\Z$, defined by \eqref{eq:kummer-block}; and $L$.
The integer-order Bessel coordinates use \eqref{eq:bessel-integer} and
their derivatives; for $\nu\in\Q\setminus\Z$ the coordinates are
$(f_{\nu,c},f'_{\nu,c},f_{-\nu,c},f'_{-\nu,c})$.
The additional $(\mathrm{B2}^E)$ restrictions belong to the arithmetic statements.
A \emph{closed} list contains every lower tower level
and each second-order block's four coordinates. It adds no (B2)
derivatives, which already lie in the exponential base.

\emph{Normalization} uses \eqref{eq:bessel-integer} and
\eqref{eq:link-order} to reduce $\nu\equiv\pm\nu_0\pmod\Z$,
$0\le\nu_0\le1/2$. The endpoint $\nu_0=1/2$ reduces to exponentials;
remaining second-order representatives have $0\le\nu_0<1/2$.
Choose one scale from $\{c,-c\}$ and retain explicit expressions for the
original coordinates. Half-order rates are $\pm c$.
For the second-order differential modules $M_i$ of these blocks,
linkage means $M_j\simeq M_i\otimes\Lambda$ over $k$ for rank-one
$\Lambda$, with a twist representative $e^{\lambda z}z^\upsilon$.
The set of linkage classes is $\mathcal C$.
Theorem~\ref{thm:linkage} decides when two blocks are linked.
A finitely generated $\Gamma\subset(\C,+)$ is \emph{saturated} for the
list when it contains the listed rates, kernel rates, every $-\eta$,
Kummer scales, half-order rates $\pm c$ and linkage twist rates.
Choose a $\Z$-basis and put
\begin{equation}\label{eq:exponential-field}
\Gamma=\bigoplus_{h=1}^d\Z\lambda_h,\quad t_h=e^{\lambda_hz},\quad
T=k(t_1,\ldots,t_d),\quad t^v=\prod t_h^{v_h}.
\end{equation}
The weight of $t^v$ is $\sum_hv_h\lambda_h$.
Put $K_{\rm coord}=k(\mathcal F)$ and $K_{\rm full}=T(\mathcal F)$;
the latter is differential for a closed list and equals $K_{\rm coord}$
when the basis exponentials are listed.

First-level coordinates are primitives with rational-exponential
derivative, including $\Ein_1(\eta z)$. For
$F'=\sum q_\mu e^{\mu z}$ define the finite-support vector
\begin{equation}\label{eq:residue-vector}
\operatorname{res}(F)=(\rho_{\mu,p}(q_\mu))_{(\mu,p)},\qquad
\rho_{\mu,p}(q)=\Res_{z=p}e^{\mu(z-p)}q(z)\,dz.
\end{equation}
Here $e_{(\mu,p)}$ is the unit vector at that pair and
$\operatorname{res}(\Phi_{\mu,p})=e^{-\mu p}e_{(\mu,p)}$. Put
\begin{equation}\label{eq:rank-data}
\mathcal V_{\rm res}=\Span_\C\{\operatorname{res}(F):F\text{ first-level}\},\quad\rho=\dim\mathcal V_{\rm res},
\quad\tau=\sum_{\eta}(h_\eta-1).
\end{equation}
Thus $\tau$ counts distinct higher pairs $(\eta,k)$, $k\ge2$.
For example, $\operatorname{res}(\Ein(\eta z))=e_{(0,0)}-e_{(-\eta,0)}$.

Let $X=(X_q)$ index every listed input slot, including repeated coordinates
such as $(\eta,k)$. For each repeated coordinate $q$, let
$\operatorname{rep}(q)$ choose one representative slot. The duplicate graph
$X_q-X_{\operatorname{rep}(q)}$ is a relation in the declared polynomial ring.
A \emph{constant coordinate} is identically $c\in\C$, with graph $X_c-c$.
Subscripts name coordinates, not values; generators may use coordinate names
as variables. \emph{Selected} coordinates mean a first-level residue basis,
all distinct higher levels and one triple per linkage class, specified in
Section~\ref{sec:main-statement}.
The dimension of a prime ideal means the Krull dimension of its quotient
ring. For \emph{Laurent saturation}, clear negative powers of the
exponential variables. Then replace the resulting ideal $J$ by
$J:m^\infty=\{f:m^kf\in J\text{ for some }k\ge0\}$, where $m$ is
the product of these variables. Elimination then gives relations among the listed coordinates.
Residue constants are additive; $\mathbf c$ denotes the connection column
of Section~\ref{sec:gamma}, and unindexed $\pi$ denotes the circle constant.

\begin{theorem}\label{thm:main}
Let $\mathcal F$ be closed and normalized as above, and let $\Gamma$
be saturated for it. All first-level primitive kernels are rational-exponential
functions with complex coefficients; higher tower coordinates are defined by
$(\mathrm{B2}^\tau)$ in \eqref{eq:ein-tower}. All Bessel orders are rational.
Bessel and Kummer scales are nonzero; Kummer first parameters are nonintegral.
Put $r=\rho+\tau+3\#\mathcal C$. Then
\begin{equation}\label{eq:main-degree}
\trdeg_TK_{\rm full}=r,\qquad\trdeg_{\C(z)}K_{\rm full}=d+r.
\end{equation}
The second degree is that of $K_{\rm coord}$ when the basis exponentials
are listed. The full prime relation ideal over $T$ has dimension
$\rho+\tau+3\#\mathcal C$. Its generators include affine
first-level residue relations and one Wronskian determinant per linkage class.
They also include linear coordinate identities for linked blocks and listed exponentials.
Equality graphs identify duplicate slots of every type.
A listed constant $c$ has relation $X_c-c$.
Half-integer Bessel coordinates have the exponential graphs of
Section~\ref{sec:main-statement}, with $\pm c\in\Gamma$.
These generators and a transcendence basis are given
explicitly in Section~\ref{sec:main-statement}. Over $\C(z)$ the full
coordinate ideal is obtained by Laurent saturation and elimination.
\end{theorem}

\begin{theorem}[Normalized linkage maps]\label{thm:linkage}
All Kummer blocks here are ${}_1F_1(a;1;\cdot)$.
After normalization, $0\leq\nu_0<1/2$. For second-order blocks, a rank-one
twist is possible precisely when
\begin{align}
 \operatorname{Bess}(c,\nu_0)&\sim\operatorname{Bess}(c',\nu'_0)
 &&\Longleftrightarrow c'^2=c^2,\quad\nu'_0=\nu_0,\label{eq:link-bb}\\
 \operatorname{Kum}(s,a)&\sim\operatorname{Kum}(s',a')
 &&\Longleftrightarrow
 \begin{cases}s'=s,\ a'-a\in\Z,\\s'=-s,\ a'+a\in\Z,\end{cases}\label{eq:link-kk}\\
 \operatorname{Bess}(c,\nu_0)&\sim\operatorname{Kum}(s,a)
 &&\Longleftrightarrow\nu_0=0,\ s=\pm2c,\ a\in\tfrac12+\Z.\label{eq:link-bk}
\end{align}
A nonresonant Bessel block is not linked to a resonant one. Every such
twist has a representative $e^{\lambda z}z^\upsilon$ after a rational gauge.
Its rate is $0$ for Bessel links
and contiguity, $\pm s$ for the Kummer transformation, and $\pm c$
for the mixed link. Its power satisfies $\upsilon\in\Z$ in resonant
classes, $\upsilon\equiv0$ or $\pm2\nu_0\pmod\Z$ otherwise. Here $\upsilon$
compares the original blocks; normalized representatives have
$\upsilon\equiv0\pmod\Z$.
Section~\ref{sec:linkage} gives explicit identities for regular solutions,
normalized entire parts (E-parts) and derivative coordinates.
Their compositions obey \eqref{eq:link-composition-convention} and
\eqref{eq:link-composed}. That section fixes the constants and orientation rules.
\end{theorem}
The equivalences underlying this statement are due to Gorelov
\cite[Lemmas~12 and~14, p.~7]{Gorelov2021}; our contribution is the explicit
normalized maps.

Gorelov's rescaled, exponentially twisted identity uses
${}_1F_1(\nu+\tfrac12;2\nu+1;\cdot)$ \cite[(3), p.~260]{Gorelov2019}.
Its second parameter is one only at $\nu=0$.
Appendix~\ref{app:algorithm} refers to Supplementary Section~\ref{sup:input-model}
for generators, maps and constants.

Use Section~\ref{sec:setup}'s $E$-coordinates and $(\mathrm{B2}^E)$ primitives,
with bounds in Appendix~\ref{app:arithmetic}. Let $X$ index coordinates.
Let $\rho^E$ be the primitive residue rank, including level-one exponential integrals. Put
$\mathcal E=\operatorname{acl}_{\C}\Qbar(e^\beta:\beta\in\Qbar)$,
where algebraic closure is taken within $\C$.
Section~\ref{sec:main-statement} selects coordinates; $y_{\mathfrak c}$ is
the representative's regular first coordinate.
Independent invertible indeterminates $\Theta=(\Theta_1,\ldots,\Theta_d)$
represent basis exponentials; put $R_\Theta=\Qbar[\Theta^{\pm1}]$.
Types 2--4 are Section~\ref{sec:main-statement}'s affine residue relations,
determinants and linkage identities.

\begin{theorem}\label{thm:values}
Let $\mathcal F$ be closed and normalized, with $\Gamma$ saturated.
All rates, nodes, scales and listed constants are algebraic.
All Bessel/Kummer parameters are rational.
The coordinates are $E$-coordinates, apart from integer-order Bessel blocks
reduced to order zero by Lemma~\ref{lem:integer}.
For $\alpha\in\Qbar^\times$, put $X(\alpha)=(q(\alpha))_{q\in\mathcal F}$
and specialize only
$z=\alpha$ in Theorem~\ref{thm:main}'s generators over
$\Qbar[z^{\pm1}][\Theta^{\pm1}]$, obtaining $\mathfrak P_\alpha$ in
$R_\Theta[X]$. Suppose the basis exponentials are listed.
Define the evaluation map
\[
 \operatorname{ev}_\alpha:R_\Theta[X]\longrightarrow\C,
 \qquad \Theta_h\longmapsto e^{\lambda_h\alpha},\quad
 X_q\longmapsto q(\alpha).
\]
Let $I_\alpha=\ker(\operatorname{ev}_\alpha)$. Then
\begin{equation}\label{eq:value-ideal}
I_\alpha=\mathfrak P_\alpha,\qquad
\trdeg_{\Qbar}\Qbar(X(\alpha))=d+\rho^E+\tau+3\#\mathcal C.
\end{equation}\equationalias{eq:intro-value-degree}
The selected coordinates and basis exponentials are independent.
In particular,
\begin{equation}\label{eq:value-nonzero}
y_{\mathfrak c}(\alpha)\ne0.
\end{equation}
For several algebraic evaluation points $\alpha_i\ne0$, the same
assertions apply to the combined list of functions $q_i(\alpha_i z)$
at $z=1$, after closure, normalization and saturation. Nodes, scales
and rates are rescaled, derivative factors are retained, and all
selection conditions are tested on this combined list.
\end{theorem}

\begin{corollary}\label{cor:value-base}
Under the hypotheses of Theorem~\ref{thm:values}, write $X_{\rm ne}$ for the variables of nonexponential coordinates.
For these coordinates,
\begin{equation}\label{eq:value-exponential-base}
\trdeg_{\mathcal E}\mathcal E(X_{\rm ne}(\alpha))
=\rho^E+\tau+3\#\mathcal C.
\end{equation}
The ideal in $\mathcal E[X_{\rm ne}]$ has the nonexponential generators
of $\mathfrak P_\alpha$, with $\Theta_h=e^{\lambda_h\alpha}$.
These include types 2--4 and the half-integer exponential graphs,
together with their derivative and order-shift graphs.
They also include all duplicate graphs and $X_c-c$ for listed algebraic constants.
Equivalently, extend the full ideal to $\mathcal E[X]$ under this
substitution, then contract to $\mathcal E[X_{\rm ne}]$.
The selected values are independent over $\mathcal E$.
At $\alpha=1$, choose first-level values with jointly independent residue vectors.
Add distinct higher levels $(\eta,k)$ and one selected triple per class.
Every finite subcollection is independent over $\mathcal E$.
The degree and ideal assertions apply to each finite closed sublist.
For mixed points $q_i(\alpha_i)$, apply the conditions to the combined
rescaled list $q_i(\alpha_i z)$.
Require independent residues, distinct higher pairs $(\eta\alpha_i,k)$
and one selected triple per class.
Conditions on the unrescaled lists alone do not suffice.
\end{corollary}

For nonclosed lists, close, normalize, adjoin basis exponentials, saturate,
specialize and eliminate auxiliary coordinates, retaining duplicate/constant graphs.
The enlarged list of values need not have the same transcendence degree as the original list.

Theorem~\ref{thm:measure} in Appendix~\ref{sec:measure} gives a linear independence estimate
for $1,e,\allowbreak e^{-1},\Ein(1)$ with evaluated constants.
\setcounter{theorem}{4}
Our $H^{-60}/12$ estimate is weaker than the classical one
and applies for $H\ge H_0$.
Supplementary Section~\ref{sup:measure-comparison} compares the classical exponent and
effectiveness and gives a cruder bound for $1\le H<H_0$.

Theorem~\ref{thm:connection} gives a multi-node, mixed-family bound over
$\mathcal E$ and the regular-value base $B$.
The one-point consequence over the full $\mathcal E$ is known
\cite[Ch.~7, \S1, Theorem~1, pp.~232--233]{Shidlovskii1989};
Supplementary Section~\ref{sup:classical-comparison} gives the finite-rate reduction and distinguishes
Mahler's announcement from Rivoal's proof.
We claim no novelty for the mixed formulation over $B$.
Section~\ref{sec:gamma} states the bound and cancellation limitations.

\subsection{Relation to earlier work}
For primitive and homogeneous antecedents see Supplementary Section~\ref{sup:classical-comparison} and Appendix~\ref{app:gorelov}
\cite{Ostrowski1946,Kolchin1973,BBH1988,Gorelov2021}.
Mahler also proved that $\pi Y_0(2)/(2J_0(2))-\gamma$ is transcendental
\cite[Theorem~3 and its consequence, pp.~169--170]{Mahler1968},
a precedent for the Bessel connection row.
Fischler--Rivoal give a general relation-ideal algorithm
\cite[Theorem~4(i)--(ii), pp.~2--3 of the arXiv version]{FischlerRivoal2019}.
Supplementary Section~\ref{sup:classical-comparison} compares the regular families,
their normalizations and the restricted inhomogeneous results.

Twisted residues, integration and $\SL_2$ factors determine the pointwise
stabilizer of the coordinate field (Lemma~\ref{lem:stabilizer}).
Integral quotients and equal dimensions then identify the full ideals.

\section{Setup}\label{sec:setup}
The data conventions preceding Theorem~\ref{thm:main} fix $k,D$,
the meromorphic domain, base point, logarithm branch and primitive
normalizations. They also define residues \eqref{eq:residue-vector},
the rate field \eqref{eq:exponential-field} and ranks \eqref{eq:rank-data}.
This section gives the coordinate dictionary used there.

\subsection{Admissible blocks}\label{subsec:blocks}
The five block types below have complex parameters $\lambda,\mu,c,s$,
with $c,s\ne0$.
\begin{description}\setlength{\itemsep}{0pt}\setlength{\parsep}{0pt}
\item[(B1)] $e^{\lambda z}$, $\lambda\in\C$.
\item[(B2)] A specified primitive of $\sum_{\mu\in M}q_\mu e^{\mu z}$,
$M$ finite, $q_\mu\in\C(z)$, with arbitrary complex additive constant.
The subclass $(\mathrm{B2}^E)$ consists of primitives with algebraic
rates, kernels $q_\mu\in\Qbar[z,z^{-1}]$, entire derivative and $F(0)=0$.
\par\noindent
The extension $(\mathrm{B2}^\tau)$ also permits finite towers. For each
$\eta\ne0$ and integer $h\ge1$, include levels $1\le k\le h$ defined by
\begin{equation}\label{eq:ein-tower}
\Ein_k(u)=\sum_{n\ge1}\frac{(-1)^{n+1}u^n}{n^kn!},\quad
D\Ein_k(\eta z)=\Ein_{k-1}(\eta z)/z\ (k\ge2).
\end{equation}
$D\Ein_1(\eta z)=(1-e^{-\eta z})/z$. Write $\Ein=\Ein_1$.
\item[(B3)] $\operatorname{Bess}(c,\nu)$, $c\ne0$, $\nu\in\Q$.
An integer Bessel order is \emph{resonant}. The local homogeneous equation has a logarithmic solution at zero. For order zero take $(I_c,V_c,S_c,W_c)$, where
\begin{equation}\label{eq:bessel-zero}
\begin{gathered}
I_c=\sum_{n\ge0}\frac{(cz/2)^{2n}}{(n!)^2},\quad
S_c=\sum_{n\ge0}\frac{\mathsf H_n(cz/2)^{2n}}{(n!)^2},\\
V_c=I_c',\ W_c=S_c',\ \mathsf H_n=\sum_{j=1}^n1/j,\ \mathsf H_0=0.
\end{gathered}
\end{equation}
As a function of its own argument, write $\mathsf S(u)=S_1(u)$; thus
$S_c(z)=\mathsf S(cz)$.
The operator $\mathscr B_{c,0}=zD^2+D-c^2z$ satisfies
$\mathscr B_{c,0}I_c=0$ and $\mathscr B_{c,0}S_c=2I_c'$. Thus $J_c=S_c-I_cL$
is a homogeneous solution. To define the other integer orders, put $R_n=a_n+b_nD$ for $n\ge0$, with
$(a_0,b_0)=(1,0)$, $(a_1,b_1)=(0,c^{-1})$ and
$(a_{n+1},b_{n+1})=(a_{n-1},b_{n-1})-2n(a_n,b_n)/(cz)$. Then
\begin{equation}\label{eq:bessel-integer}
I_{n,c}=R_nI_c=I_n(cz),\quad P_{n,c}=R_nS_c-b_nI_c/z,\quad
R_nJ_c=P_{n,c}-I_{n,c}L,
\end{equation}
with their derivatives, define the integer block. Use $I_{-n}=I_n$ and $P_{-n,c}=P_{n,c}$;
$P_{n,c}$ may have a pole at zero, e.g. $P_{1,c}=(W_c-I_c/z)/c$.
For every rational $\nu$ that is not a negative integer put
\begin{equation}\label{eq:bessel-nonresonant}
f_{\nu,c}=\sum_{n\ge0}\frac{(cz/2)^{2n}}{n!(\nu+1)_n},\qquad
\mathscr B_{c,\nu}=zD^2+(2\nu+1)D-c^2z.
\end{equation}
For every $\nu\in\Q\setminus\Z$, use the coordinate quadruple
$(f_{\nu,c},f'_{\nu,c},f_{-\nu,c},f'_{-\nu,c})$.
An order $\nu\notin\tfrac12\Z$ is \emph{nonresonant}: its two local
monodromy eigenvalues at zero are distinct.
The two solutions in this case are
$f_{\nu,c}=\Gamma(\nu+1)(cz/2)^{-\nu}I_\nu(cz)$ and
$z^{-2\nu}f_{-\nu,c}$. Half-integer orders reduce to exponentials.
The starting identities are
\[ f_{1/2,c}=\frac{\sinh(cz)}{cz},\qquad f_{-1/2,c}=\cosh(cz). \]
\item[(B4)] $\operatorname{Kum}(s,a)$, $s\ne0$, $a\notin\Z$, is resonant, with a
logarithmic homogeneous solution at zero. Its coordinates are
\[ (F,F',\Sigma,\Sigma')=(f_a(sz),D[f_a(sz)],S_a(sz),D[S_a(sz)]). \]
Here $f_a={}_1F_1(a;1;\cdot)$; the second Kummer parameter is fixed to one.
\begin{equation}\label{eq:kummer-block}
f_a(u)=\sum_{n\ge0}\frac{(a)_nu^n}{(n!)^2},\quad
S_a(u)=\sum_{n\ge0}\frac{(a)_nh_nu^n}{(n!)^2},\quad
h_n=2\mathsf H_n-\sum_{j<n}\frac1{j+a}.
\end{equation}
Here $(a)_0=1$ and $(a)_n=a(a+1)\cdots(a+n-1)$.
The operator $\mathscr K_{s,a}=zD^2+(1-sz)D-as$ satisfies
$\mathscr K_{s,a}F=0$ and $\mathscr K_{s,a}\Sigma=2F'-sF$. Thus $J=\Sigma-FL$ is
homogeneous. In the arithmetic statements we require $a\in\Q\setminus\Z$.
\item[(B5)] $L$, the primitive of $1/z$.
\end{description}
The tower at $\eta\ne0$ has depth $h_\eta$ in \eqref{eq:ein-tower};
closure is defined before Theorem~\ref{thm:main}.

The following are $E$-coordinates: listed exponentials with algebraic
rates, $(\mathrm{B2}^E)$ primitives, and tower levels $\Ein_k(\eta z)$
with $\eta\in\Qbar^\times$.
For Bessel blocks, require algebraic $c$ and rational order.
Include noninteger-order coordinates and the order-zero coordinates
$(I_c,V_c,S_c,W_c)$.
Other integer-order blocks contribute only the regular coordinates
$I_{n,c},I'_{n,c}$.
Kummer coordinates have algebraic $s$ and rational nonintegral $a$;
listed algebraic constants are also allowed. The integer-order
companions $P_{n,c},P'_{n,c}$ for $n\ge1$ are excluded and are handled
by order-zero reduction. The functions $L$, $G_{\eta,k}$ and
the logarithmic solutions $J$ are not $E$-coordinates.

Normalization, saturation, coordinate variables, duplicate/constant
graphs and Laurent elimination are defined together before
Theorem~\ref{thm:main}. Lemma~\ref{lem:exponential-log} proves
exponential independence.

\section{The degree and the relation ideal}\label{sec:main-statement}
Put $r=\rho+\tau+3\#\mathcal C$, as in Theorem~\ref{thm:main}.
A transcendence basis contains a first-level residue basis.
It also contains all distinct levels $(\eta,k)$ with $k\ge2$.
For each class it contains one triple:
$(I_c,V_c,S_c)$, $(F_{s,a},F'_{s,a},\Sigma_{s,a})$, or
$(f_{\nu,c},f'_{\nu,c},f_{-\nu,c})$. It stays independent over $T(L)$
when $e_{(0,0)}\notin\mathcal V_{\rm res}$. Then $L$ is transcendental over
$K_{\rm full}$. Otherwise $L\in K_{\rm full}$.

Use the coordinate variables defined before Theorem~\ref{thm:main}.
The full prime ideal $\mathfrak p\subset T[X]$ has dimension $r$.
Its generators are:
\begin{enumerate}\setlength{\itemsep}{0pt}\setlength{\parsep}{0pt}
\item $X_e-t^{v_e}$ for listed exponentials of rate $\sum(v_e)_h\lambda_h$;
\item for nonbasis first-level $F_j$,
\begin{equation}\label{eq:residue-relations}
F_j-\sum_i\kappa_{ji}F_i-r_j-c_j,\quad
\operatorname{res}(F_j)=\sum_i\kappa_{ji}\operatorname{res}(F_i),\quad r_j\in\C(z)[t^{\pm1}],\ c_j\in\C;
\end{equation}
\item one representative determinant,
\begin{equation}\label{eq:determinant-relations}
\begin{array}{ll}
\operatorname{Bess}(c,0):&z(IW-VS)-I^2+1,\\
\operatorname{Kum}(s,a):&z(F\Sigma'-F'\Sigma)-F^2+e^{sz},\\
\operatorname{Bess}(c,\nu):&z(f_\nu f'_{-\nu}-f'_\nu f_{-\nu})-2\nu f_\nu f_{-\nu}+2\nu;
\end{array}
\end{equation}
\item every nonrepresentative coordinate minus its $T$-linear representative
expression from Theorem~\ref{thm:linkage} and derivatives. The last determinant
uses its own nonresonant representative order. \end{enumerate}
Include $X_q-X_{\operatorname{rep}(q)}$ for every duplicate slot of any coordinate type.
For each listed constant, include $X_c-c$.
These monic graphs eliminate duplicate and constant slots.
The degree is unchanged. Only distinct higher levels remain free.
The data definitions preceding Theorem~\ref{thm:main} define $\operatorname{rep}$.
For half-integer orders, add the exponential graphs ($\pm c\in\Gamma$)
\[
 X_{f_{1/2,c}}-\frac{\sinh(cz)}{cz},\qquad
 X_{f_{-1/2,c}}-\cosh(cz),
\]
and add their derivatives and order shifts from \eqref{eq:link-order}.
Write $\sinh,\cosh$ as Laurent polynomials in the basis exponentials.

For the Laurent form, replace $t$ by $\Theta$ in all generators.
Work over $\mathscr A_{\exp}=\C(z)[\Theta^{\pm1}]$.
The generated prime ideal is exactly
$\mathfrak P=\ker(\mathscr A_{\exp}[X]\to K_{\rm full})$ under $\Theta_h\mapsto t_h$,
of dimension $d+r$. \mbox{Its} contraction to $\C(z)[X]$ is the full coordinate ideal.
Laurent saturation and elimination compute it, including exponential toric relations.
To replace $\Theta$-monomials by listed exponentials, first clear negative powers.
Then saturate by their product. Let $k_1$ be any field containing
the data, listed constants and fundamental $c_j$.
The same generators give $\mathfrak P\cap k_1(z)[\Theta^{\pm1}][X]$.
This ideal is prime of dimension $d+r$.
The corresponding degree is $\trdeg_{k_1(z)}k_1(z)(t,\mathcal F)=d+r$.
For the normalization-compatible input of Appendix~\ref{app:algorithm},
the fundamental constants $c_j$ belong to $k_1$.
Under Theorem~\ref{thm:values}, first-level coordinates vanish at zero.
The generators lie in
$\Qbar[z^{\pm1}][\Theta^{\pm1}][X]$.


\section{Proof of Theorem~\ref{thm:main}}\label{sec:proof}
\subsection{Primitives and towers}
Put $F_0=\C(z)$, $F_1=\C(z,L)$, $R_i=F_i[t^{\pm1}]$. On $F_1$, $D=\partial_z+z^{-1}\partial_L$.
\begin{lemma}\label{lem:exponential-log}
The $t_h$ are algebraically independent over $F_1$, and $L$ is
transcendental over $T$. For $q\in F_i^\times$, $i=0,1$,
$Dq/q\in\C$ forces $q\in\C^\times$.
\end{lemma}
\begin{proof}
Suppose distinct exponentials satisfy a rational linear relation with the
fewest terms. Divide by one term, making its coefficient $1$, and
differentiate. The resulting relation has fewer terms. Thus all its
coefficients vanish. A remaining coefficient would give
$q'/q=-\mu\ne0$ for rational $q$. This is impossible, since a rational
logarithmic derivative tends to zero at infinity. Distinct monomials in
the $t_h$ have distinct weights. They are therefore linearly independent
over $F_0$. This proves algebraic independence of the $t_h$.

A polynomial relation of least positive degree $n$ for $L$ over $T$
can be continued around zero. Its coefficients return to their original
values, while $L$ changes to $L+2\pi i$. Subtraction gives a relation of
degree $n-1$, with leading coefficient $2\pi in$ times the former leading
coefficient. This is a contradiction. Thus $L$ is transcendental over
$T$, and the $t_h$ remain independent over $F_1$.

For $p=L^N+c_{N-1}L^{N-1}+\cdots$, $Dp$ has degree $N-1$ and coefficient
$c'_{N-1}+N/z\ne0$: a rational derivative has zero residue at zero.
For monic irreducible $p$ of positive degree this gives $p\nmid Dp$.
Factor $q\in F_1^\times$ over $\C(z)[L]$. Every nonzero factor exponent
then contributes an uncancelled pole to $Dq/q$. If $Dq/q$ is constant,
no such factor occurs. Thus $q\in F_0^\times$. Its logarithmic derivative
tends to zero at infinity. It is therefore zero, and $q\in\C^\times$.
This also proves the assertion for $i=0$.
\end{proof}

\begin{lemma}\label{lem:reduction}
A positive-degree monic irreducible $p\in\C(z)[L]$ satisfies
$p\nmid Dp\ne0$. No nonunit irreducible $p\in R_i$, $i=0,1$, divides
$Dp$. Consequently $y\in\operatorname{Frac}R_i$ and $Dy\in R_i$
imply $y\in R_i$.
\end{lemma}
\begin{proof}
The first assertion was proved above. Suppose $Dp=hp$ in $R_i$.
Differentiation does not enlarge the monomial support of $p$.
If $Dp\ne0$, the Newton polytope of $hp$ is the sum of those of $h$
and $p$. Its containment in the polytope of $p$ forces $h\in F_i$.
The same conclusion holds with $h=0$ when $Dp=0$.
For two nonzero coefficients $a_v,a_w$ of distinct weights
$\omega_v,\omega_w$, coefficient comparison now gives
\[
 \frac{D(a_v/a_w)}{a_v/a_w}=\omega_w-\omega_v\ne0.
\]
Lemma~\ref{lem:exponential-log} excludes this. Thus $p$ is a monomial,
which is a unit of $R_i$.

Let a denominator factor of $y$ have order $s>0$ at $p$.
In the localization at $(p)$ write $y=wp^{-s}$, with $w$ a unit. Then
\[
 D(wp^{-s})=\frac{pDw-swDp}{p^{s+1}}.
\]
The numerator is nonzero modulo $p$, because $p\nmid Dp$.
Thus $Dy$ has valuation $-s-1$ at $p$. This contradicts $Dy\in R_i$,
and removes every denominator factor.
\end{proof}

\begin{lemma}\label{lem:twisted-obstruction}
For $\mu\ne0$, $g\in\C(z)$ and nonconstant $P\in\C[L]$, the equation
$Du+\mu u=P(L)/z+g$ has no solution in $F_1$.
\end{lemma}
\begin{proof}
The denominator argument of Lemma~\ref{lem:reduction} also applies to $(D+\mu)u$. Take an irreducible $L$-denominator of order $s>0$. Then $Du$ has pole order $s+1$. In contrast, $\mu u$ has order $s$.
Hence $u\in\C(z)[L]$. Write $u=\sum_{j=0}^nu_jL^j$ and
$P=\sum_{j=0}^mp_jL^j$, where $m\ge1$ and $p_m\ne0$.
If $n>m$, its highest coefficient gives $u_n'+\mu u_n=0$,
excluded by Lemma~\ref{lem:exponential-log}. If $n<m$, the coefficient
of $L^m$ cannot match. Hence $n=m$, and
\[
 u_n'+\mu u_n=p_n/z,\qquad p_n\ne0.
\]
A finite pole of $u_n$ would increase order under differentiation,
including at zero, where the right side has only a simple pole.
Thus $u_n$ is polynomial, whose image under $D+\mu$ cannot equal $p_n/z$.
\end{proof}

\begin{lemma}[Residues and Ostrowski's trace argument]\label{lem:primitive}
For $F_j'=\sum q_{j,\mu}e^{\mu z}$ with rates in $\Gamma$, the
field generated over $T$ by these meromorphic primitives has constant
field $\C$: it introduces no new constants on the fixed domain above
Theorem~\ref{thm:main}. Then
\begin{equation}\label{eq:core-residue-criterion}
\sum b_jF_j\in T\iff\sum b_j\operatorname{res}(F_j)=0\quad(b_j\in\C).
\end{equation}
The combination then equals $r+c$, $r\in R_0$, $c\in\C$.
A residue basis is independent over $T$ and its affine relations generate
the full ideal.
\end{lemma}
\begin{proof}
For rational $q$,
\begin{equation}\label{eq:core-rational-residue}
(D+\mu)r=q\ (r\in\C(z))\iff\rho_{\mu,p}(q)=0\quad(p\text{ finite}).
\end{equation}
Necessity follows by differentiating $e^{\mu(z-p)}r$. For sufficiency,
subtract images of $-c(z-p)^{1-m}/(m-1)$ for $m\ge2$ successively. The
remaining simple-pole coefficients are the zero residues. On polynomials,
$D+\mu$ is bijective if $\mu\ne0$. If $\mu=0$, take a polynomial primitive.
Lemma~\ref{lem:reduction} reduces combinations in $T$ to Laurent form. Weight comparison proves \eqref{eq:core-residue-criterion}.

We use Ostrowski's trace argument (F3). Let dependent same-constant
primitives over $F_{\rm b}$ have a maximal algebraically independent
subfamily $y_1,\ldots,y_q$.
Put $E=F_{\rm b}(y_1,\ldots,y_q)$ and choose a remaining primitive $y$.
It is algebraic over $E$. Derivation commutes with trace. Thus $v=\operatorname{Tr}_{E(y)/E}(y)/[E(y):E]$ satisfies $Dv=Dy\in F_{\rm b}$.
Thus $y-v\in\C$. The partial $\partial_i=\partial/\partial y_i$
fixing $F_{\rm b}$ commutes with $D$, since each $Dy_j$ belongs to $F_{\rm b}$.
It follows that $D(\partial_iv)=\partial_i(Dv)=0$. Thus $\partial_iv=c_i\in\C$. A rational function with all these partials
zero lies in $F_{\rm b}$. Therefore $v-\sum c_iy_i\in F_{\rm b}$, giving a nonzero
constant linear combination of the primitives in $F_{\rm b}$. Apply this with
$F_{\rm b}=T$ and use \eqref{eq:core-residue-criterion}.
The affine quotient is the polynomial ring on a residue basis. In particular,
\begin{equation}\label{eq:core-reference}
F=\sum_{(\mu,p)}e^{\mu p}\operatorname{res}(F)_{(\mu,p)}\Phi_{\mu,p}+r_F+c_F.
\end{equation}
\end{proof}

Put
\begin{equation}\label{eq:core-G}
G_{\eta,k}=L^k/k!-\Ein_k(\eta z),\quad
DG_{\eta,1}=e^{-\eta z}/z,\quad DG_{\eta,k}=G_{\eta,k-1}/z.
\end{equation}
For this proof choose $\Phi_{-\eta,0}=G_{\eta,1}$.
This changes only the reference constants. The specified coordinates
and their additive constants retain their original values.
\begin{lemma}\label{lem:tower}
For distinct nonzero $\eta_j$, $h_j\ge1$, and a finite residue support
$\mathfrak S$ containing $(0,0),(-\eta_j,0)$ with rates in $\Gamma$, the
generators of
$\mathcal N_h=T(\Phi_\sigma:\sigma\in\mathfrak S;
G_{\eta_j,k}:2\le k\le h_j)$ are independent. This meromorphic field
has constant field $\C$, so it introduces no new constants on the fixed
domain above Theorem~\ref{thm:main}. For $P_j\in\C[L]$, not all
zero,
\begin{equation}\label{eq:core-NP}
\sum_jP_j(L)G_{\eta_j,h_j}/z\notin D\mathcal N_h.
\end{equation}
The solvable coordinates have degree $\rho+\tau$ and only their first-level
residue relations.
\end{lemma}
\begin{proof}
A detailed proof is given in Appendix~\ref{app:detailed-proofs}.
Induct on $\sum h_j$. At each depth first prove independence, then
\eqref{eq:core-NP}. When all $h_j=1$, the residue vectors of the
$\Phi_\sigma$ are nonzero multiples of distinct coordinate vectors.
Lemma~\ref{lem:primitive} gives independence, with $\Phi_{0,0}=L$.
The other coordinates therefore form a rational function field over $T(L)$.

Suppose $Dy$ equals the sum in \eqref{eq:core-NP} at this depth.
For $\sigma\ne(0,0)$ take the partial in $\Phi_\sigma$, fixing $T(L)$.
It commutes with $D$, because $D\Phi_\sigma\in T$.
For $\sigma=(-\eta_j,0)$ this gives
\[
 D(\partial_{-\eta_j,0}y)=P_j(L)/z,
 \qquad \partial_{-\eta_j,0}y=Q_j(L),\quad Q_j'=P_j.
\]
Absorb constants in $Q_j$; every other partial is constant.
Subtract $Q_j\Phi_{-\eta_j,0}$ and the other $c_\sigma\Phi_\sigma$.
All partials vanish, leaving
$y_0\in T(L)$ with
\[
 Dy_0=-\sum_jQ_j(L)e^{-\eta_jz}/z-\sum_{\sigma}c_\sigma\phi_\sigma.
\]
The second sum contains the other residue coordinates.
Lemma~\ref{lem:reduction} puts $y_0$ in $R_1$.
Its coefficient of weight $-\eta_j$ satisfies
$Du-\eta_ju=-Q_j(L)/z+g$, with $g\in\C(z)$.
Some $P_j\ne0$ makes $Q_j$ nonconstant. Distinct nodes have distinct
weights; the other residue terms contribute only rational $g$.
Lemma~\ref{lem:twisted-obstruction} excludes this equation.

Now let some $h_j\ge2$, and put $\widetilde h_j=h_j-1$,
$\widetilde h_i=h_i$ for $i\ne j$.
The generators of $\mathcal N_{\widetilde h}$ are independent by induction.
Adjoin $Y=G_{\eta_j,h_j}$, with
$DY=G_{\eta_j,\widetilde h_j}/z\in\mathcal N_{\widetilde h}$.
If $Y$ were algebraic over this field, its normalized trace would be
$v\in\mathcal N_{\widetilde h}$ with $Dv=DY$, by
Lemma~\ref{lem:primitive}: derivation commutes with trace, and $Y-v\in\C$.
Formula~\eqref{eq:core-NP} at
$\widetilde h$, with coefficient $1$ for $j$ and $0$ otherwise,
excludes this. Thus adjoining $Y$ preserves independence.

To prove \eqref{eq:core-NP} at $h$, put $J=\{j:h_j\ge2\}$.
Put $h''_j=h_j-1$ for $j\in J$, and $h''_j=1$ otherwise.
The independence just proved identifies
$\mathcal N_h=\mathcal N_{h''}(Y_j:j\in J)$ as a rational function field,
where $Y_j=G_{\eta_j,h_j}$ and $DY_j=G_{\eta_j,h''_j}/z$.
The partials $\partial_j=\partial/\partial Y_j$ fixing
$\mathcal N_{h''}$ commute with $D$ because every $DY_j$ lies in that field.
If $y$ were the alleged primitive, then
$D(\partial_jy)=P_j(L)/z=DQ_j(L)$ for a polynomial $Q_j$ with
$Q_j'=P_j$. Since all constants are $\C$, an additive constant makes
$\partial_jy=Q_j(L)$. Hence
$y_0=y-\sum_{j\in J}Q_jY_j$ has every $\partial_jy_0=0$.
A rational function with these partials zero lies in the coefficient
field $\mathcal N_{h''}$. Differentiation gives
\[
\sum_{h_j=1}P_j\Phi_{-\eta_j,0}/z-
\sum_{h_j\ge2}Q_jG_{\eta_j,h_j-1}/z.
\]
If a nonzero $P_j$ has $h_j=1$, it survives in the first sum.
Otherwise some $Q_j$ in the second sum is nonconstant.
Disjoint node indices prevent cancellation. This contradicts
\eqref{eq:core-NP} at the strictly smaller height $h''$.

Finally choose $\mathfrak S$ to contain every first-level support.
Over $T(\Phi_\sigma:\sigma\in\mathfrak S)$, replacing each higher
$G_{\eta,k}$ by $\Ein_k(\eta z)=L^k/k!-G_{\eta,k}$ preserves independence.
A residue basis has degree $\rho$, and these higher coordinates add
$\tau$. This triangular replacement has an inverse over the same field;
reference changes affect only constants, preserving the specified primitives.
Every remaining first-level coordinate is its affine combination
by Lemma~\ref{lem:primitive}. This proves the degree and ideal assertions.
\end{proof}

\subsection{Groups and twists}
Let $\ell$ be the lcm of nonresonant $2\nu_0$ denominators ($\ell=1$ if
none), and $T'=T(z^{1/\ell})$. All constants remain $\C$.
\begin{lemma}\label{lem:groups}
The groups are $\SL_2$ for order-zero Bessel, have identity component
$\SL_2$ for nonresonant Bessel, and are $\GL_2$ for nonintegral Kummer.
\end{lemma}
\begin{proof}
Bessel has algebraic Wronskian. Thus $G^0\subset\SL_2$.
The algebraic gauge $z^\nu$ and $x=icz$ reduce to (F6), which excludes
a nonzero Liouvillian solution. If $G^0$ were solvable, the Liouvillian criterion in (F5) would make
the Picard--Vessiot (PV) field Liouvillian \cite[Theorem~1.43, pp.~33--34]{VanDerPutSinger2002Draft}. Every solution would then be Liouvillian. Thus $G^0$ is
not solvable. By (F15), every proper connected subgroup of $\SL_2$
is solvable. Thus $G^0=\SL_2$. In order zero its Wronskian
is rational. Hence $G=\SL_2$.
For Kummer rescale $s=1$. A rational invariant line would give
\begin{equation}\label{eq:core-Riccati}
r'+r^2+(z^{-1}-1)r-a/z=0.
\end{equation}
Its nonzero finite poles have residue $1$. Zero cannot be a pole.
At infinity $r=\epsilon+m/z+O(z^{-2})$, $\epsilon=0$ or $1$, with
$m\ge0$ the number of finite poles. Substitution gives $m=-a$ or $a-1$,
contrary to $a\notin\Z$. Thus $G$ is irreducible.
At zero the indicial equation is $r^2=0$, with resonant exponents $(0,0)$.
The local solution $J=S_a-f_aL$, with $f_a(0)=1$, changes to
$J-2\pi i f_a$ around zero. This gives a nonidentity unipotent in $G$. The Zariski closure of its powers is connected by (F15). Thus it lies in
$G^0$. If $G^0$ fixed a line, this
would be its unique fixed line and normality would make it $G$-stable.
So $G^0$ is irreducible. Its unipotent radical is normal
\cite[\S11.21, p.~157]{Borel1991}.
Its representation image is unipotent \cite[\S4.5, p.~85]{Borel1991}
and has a nonzero fixed vector \cite[\S17.5, Theorem, p.~112]{Humphreys1975}.
Normality makes this fixed space $G^0$-stable by conjugation.
Irreducibility makes it the full solution space, so the faithful radical
action is trivial. By (F15), $G^0\supset\SL_2$.
For the original scale $s$, the determinant is $e^{sz}/z$.
Its transcendence makes the determinant character of $G$ surjective
onto $\mathbb G_m$. Together with $G^0\supset\SL_2$, this gives $G=\GL_2$.
\end{proof}
The formal comparison used below is Lemma~\ref{lem:formal}, proved in
Appendix~\ref{app:foundations}.
\begin{lemma}\label{lem:goursat}
A subdirect Lie algebra in $\mathfrak{sl}_2\oplus\mathfrak{sl}_2$ is
either the full sum or the graph of an inner automorphism.
\end{lemma}
\begin{proof}
Its intersections with the two factors are ideals in those factors.
If either intersection is nonzero, simplicity and surjectivity give
the full sum. Otherwise both projections are isomorphisms. Thus the
subalgebra is the graph of an automorphism. The elementary proof after
(F15) shows that every automorphism of $\mathfrak{sl}_2$ is inner.
\end{proof}

\begin{lemma}\label{lem:twist}
The joint projective Lie algebra for two blocks is proper subdirect
exactly when their modules differ by a rank-one twist over $k$.
The twist may be written $e^{\lambda z}z^\upsilon$.
\end{lemma}
\begin{proof}
Let $G$ be the joint PV group, with solution representations $\pi_i,
\pi_j$ on $V_i,V_j$. Lemma~\ref{lem:groups} and (F15) make the
projective Lie algebra subdirect. If it is proper,
Lemma~\ref{lem:goursat} writes it as the graph of $\operatorname{Ad}A$.
The projective image of $G^0$ is closed and connected. It has the same
Lie algebra as $\{(x,AxA^{-1}):x\in\PGL_2\}$.
By (F15) the two connected groups are equal. Hence, on $G^0$,
$\pi_j=\chi A\pi_iA^{-1}$ for a scalar character $\chi$.

The representation on $\operatorname{Hom}(V_i,V_j)$ is therefore
$(1\oplus\operatorname{ad})\otimes\chi$.
Its unique one-dimensional invariant subspace is $\C A$. The traceless summand is the irreducible adjoint representation of $\mathfrak{sl}_2$.
The full group normalizes $G^0$. Thus it preserves this unique line.
By (F5) the line descends to a rank-one differential submodule over $k$.
Write $Y_i'=A_iY_i$, $Y_j'=A_jY_j$. A rational basis of that line is
an invertible matrix $B_0$ satisfying
\[
 B_0'+B_0A_i-A_jB_0=rB_0,\qquad
 B_0=wY_jAY_i^{-1},\qquad w'/w=r.
\]
Invertibility follows from the displayed expression. Multiplication by
a scalar solution of $v'/v=-r$ turns $B_0$ into an intertwiner.
Thus the twist coefficient is $r_t=-r$.
Conversely, a rank-one twist gives the same projective representation. Thus its joint projective Lie algebra is a graph.

For the shape of the twist, put $w_i=\operatorname{tr}A_i
=\sigma_i+e_i/z$ and $f=\det B_0$. Taking traces gives
\begin{equation}\label{eq:core-twist-shape}
2r_t=w_j-w_i-f'/f,\qquad f\in k^\times.
\end{equation}
Both block systems are ordinary at every finite $p\ne0$.
Their trivial local monodromy forces the twist monodromy to be trivial,
as in Lemma~\ref{lem:formal}. Formula~\eqref{eq:core-twist-shape}
therefore excludes odd $\operatorname{ord}_p f$.
Factoring $f$ gives $f=\kappa z^ng^2$, with $g\in k^\times$.
The same formula now reads
\[
 r_t=\lambda+\upsilon/z-g'/g.
\]
A rational gauge removes $g^{-1}$, leaving the scalar solution
$e^{\lambda z}z^\upsilon$.
\end{proof}
Theorem~\ref{thm:linkage} now gives exactly the classes. Choose representative
matrices $Y_{\mathfrak c}$ with columns $(y,J=\Sigma-yL)$ or
$(f_\nu,z^{-2\nu}f_{-\nu})$ and their derivatives. The maps give
\begin{equation}\label{eq:core-block-transport}
Y_i=B_iY_{\mathfrak c}w_iC_i,\quad B_i\in\GL_2(k),\ C_i\in\GL_2(\C),\quad
w_i=e^{\lambda_i z}z^{\upsilon_i}\in T'.
\end{equation}
The saturation of $\Gamma$ includes each rate $\lambda_i$.
The choice of $\ell$ includes the required algebraic powers. Thus every
scalar $w_i$ belongs to $T'$.

\subsection{The full joint group}
Let $\mathcal U/T'$ be generated by $L$, the required $\Phi_\sigma$ and
$G_{\eta,k}$. Primitive matrices are
\[
 \begin{pmatrix}1&0\\\Phi&1\end{pmatrix}.
\]
A node matrix has first column $(1,G_1,\ldots,G_h)^t$ and
\begin{equation}\label{eq:core-unipotent-matrix}
Z_{kj}=L^{k-j}/(k-j)!\ (1\le j\le k\le h),\qquad Z_{kj}=0\ (k<j).
\end{equation}
Its equations are $x_0'=0$, $x_1'=e^{-\eta z}x_0/z$, $x_k'=x_{k-1}/z$.
Together with the logarithm matrix these generate $\mathcal U$, a PV field
with connected unipotent group $\mathsf G_{u,\mathcal U}$. Put $Q=k(Y_i)$,
$G_Q=\Gal(Q/k)$ and $\Omega=\mathcal U Q$. Over $k$, add the exponential
and algebraic rank-one systems, and use first columns
$(e^{\mu z},\Phi_{\mu,p})^t$ and $(e^{-\eta z},G_1,\ldots,G_h)^t$.
They exhibit $\Omega/k$ as PV.

We record explicitly the base change from $k$ to $T'$.
The latter is PV for the exponential and algebraic rank-one systems;
its group is a subgroup of $\mathbb G_m^d\times\mu_\ell$, hence abelian.
For a block PV field $\mathcal P_i/k$, restriction identifies its group over
$T'$ with $\Gal(\mathcal P_i/(\mathcal P_i\cap T'))$. This is normal in the block group
and has abelian quotient. It therefore contains the derived group of
the identity component. For Bessel this is $\SL_2$; for Kummer,
every abelian quotient of $\GL_2$ factors through the determinant,
whose kernel is $\SL_2$. Each block Wronskian lies in $T'$:
it is a constant multiple of $1/z$, $z^{-2\nu-1}$ or $e^{sz}/z$.
Thus the base-changed group is also contained in $\SL_2$.
In particular, its identity component remains $\SL_2$.
\begin{proposition}\label{prop:product}
Consider the differential Galois group of the compositum $\Omega$ over $T'$.
It is a direct product of the unipotent primitive group and the
second-order factors. Each linkage class supplies one $\SL_2$ factor.
More precisely,
\begin{equation}\label{eq:core-product}
\Gal(\Omega/T')=\mathsf G_{u,\mathcal U}\times\prod_{\mathfrak c}\SL_2^{(\mathfrak c)}.
\end{equation}
Each class factor acts through the conjugations $C_i$ in
\eqref{eq:core-block-transport}.
\end{proposition}
\begin{proof}
A detailed proof is given in Appendix~\ref{app:detailed-proofs}.
Write $\sigma(Y_i)=Y_i\pi_i(\sigma)$ and put
$N_{\rm ss}=\Gal(QT'/T')$. Restriction embeds $N_{\rm ss}$ in $G_Q$,
because the entries of the $Y_i$ generate $QT'$ over $T'$. Let $\psi$
be the projection from $\prod_i\GL_2$ to $\prod_i\PGL_2$. Its kernel
consists of tuples of scalar matrices and is central.
Base change supplies individual $\SL_2$ projections; class independence remains to be proved.

Every differential $k$-automorphism of $QT'$ or $\Omega$ preserves $T'$. It sends
$t_h\mapsto c_ht_h$ and $z^{1/\ell}\mapsto\zeta z^{1/\ell}$,
where $c_h\in\C^\times$ and $\zeta^\ell=1$. This follows by taking
derivatives of the ratios. Two such differential $k$-automorphisms commute on $T'$.
The commutator of these differential $k$-automorphisms therefore fixes $T'$ pointwise.

In \eqref{eq:core-block-transport}, the scalar $w_i$ also belongs to $Q$:
the identity
$w_i\mathbf1=Y_{\mathfrak c}^{-1}B_i^{-1}Y_iC_i^{-1}$ proves this.
Its logarithmic derivative lies in $k$. Hence
\begin{equation}\label{eq:core-character}
 \pi_i(\sigma)=\chi_i(\sigma)C_i^{-1}\pi_{\mathfrak c}(\sigma)C_i,
 \qquad \chi_i=\sigma(w_i)/w_i,\qquad \chi_i|_{N_{\rm ss}}=1.
\end{equation}
Let $P$ be the subgroup with representative coordinates in $\SL_2$
and all other coordinates their conjugates by $C_i$. Thus
$P\simeq\prod_{\mathfrak c}\SL_2$, and
$\bar P:=\psi(P)\simeq\prod_{\mathfrak c}\PGL_2$. Both groups are closed
and connected, of dimension $3\#\mathcal C$.
Equation~\eqref{eq:core-character} gives $\psi(G_Q)\subset\bar P$.

Restriction from $G_Q$ to the group of each block is onto. The field $Q$ is the compositum of their PV fields. By (F15), the identity components
and their Lie algebras also map onto. Thus the projective Lie algebra
$\mathfrak h$ projects onto each $\mathfrak{sl}_2$.
Surjectivity makes each radical image a solvable ideal in a simple factor,
hence zero. The joint inclusion is injective, so the radical vanishes
and $\mathfrak h$ is semisimple.

Decompose $\mathfrak h$ into simple summands. A nonzero projection of one summand is an ideal in $\mathfrak{sl}_2$. Hence it is all of it. The projection is injective by simplicity. Two summands cannot both
project nontrivially to one factor, since their images commute.
Each coordinate therefore selects exactly one summand, isomorphically.
Every summand occurs in some coordinate, since the joint inclusion is faithful.
Coordinates selecting the same summand form a diagonal graph, as in
Lemma~\ref{lem:goursat}. By Lemma~\ref{lem:twist}, two such blocks are
linked. Different classes select different summands, giving
$\dim\mathfrak h\ge3\#\mathcal C$. The reverse inequality follows
from $\psi(G_Q)\subset\bar P$. Equality of dimensions in the connected
group $\bar P$, by (F15), gives
\begin{equation}\label{eq:core-projective-group}
 \psi(G_Q^0)=\psi(G_Q)=\bar P.
\end{equation}

We prove $N_{\rm ss}=P$ as groups, retaining finite components.
The Wronskians listed in the base-change paragraph belong to $T'$.
Thus $N_{\rm ss}$ has determinant one in every block, and
\eqref{eq:core-character}
gives $N_{\rm ss}\subset P$.

The extension $QT'/k$ is PV for the block and rank-one systems.
By (F5), restriction $\Gal(QT'/k)\to G_Q$ is onto.
Thus any $g,h\in G_Q$ have lifts $\widetilde g,\widetilde h$.
Their commutator fixes $T'$ because their restrictions to $T'$ commute.
It therefore belongs to $N_{\rm ss}$, and its restriction to $Q$ is
$[g,h]$. Hence $[g,h]\in N_{\rm ss}$ under the restriction embedding.
This uses lifts in the full PV group, without discarding finite components.

For $x,y\in P$, \eqref{eq:core-projective-group} supplies $g,h\in G_Q$ with
$\psi(g)=\psi(x)$ and $\psi(h)=\psi(y)$.
The factors $z_1=x^{-1}g$, $z_2=y^{-1}h$ belong to the central kernel
of $\psi$. Thus $g=xz_1$, $h=yz_2$ and central factors cancel in
$[x,y]=[g,h]\in N_{\rm ss}$.
The group $\SL_2(\C)$ is generated by upper and lower elementary
unipotents. Each upper elementary matrix is a commutator:
\[
 \begin{pmatrix}1&b\\0&1\end{pmatrix}
 =\left[\begin{pmatrix}t&0\\0&t^{-1}\end{pmatrix},
         \begin{pmatrix}1&b/(t^2-1)\\0&1\end{pmatrix}\right]
 \qquad(t^2\ne1).
\]
The lower case is identical. Thus $P$ is perfect. This proves $P\subset
N_{\rm ss}$. Both inclusions concern the full groups. Thus $N_{\rm ss}=P$ also determines every finite component.

Restriction embeds $\Gal(\Omega/T')$ as a subdirect subgroup
$H\subset \mathsf G_{u,\mathcal U}\times P$. Put
$A_{\rm u}=\{u:(u,1)\in H\}$ and $A_P=\{g:(1,g)\in H\}$.
Surjectivity of the two projections permits conjugation by lifts,
showing that $A_{\rm u}\triangleleft \mathsf G_{u,\mathcal U}$ and $A_P\triangleleft P$.
The assignment $uA_{\rm u}\mapsto gA_P$, for $(u,g)\in H$, is well defined:
two lifts of $u$ differ by $(1,g^{-1}g')\in H$. Changing $u$ by $A_{\rm u}$
does not change the target coset. It is onto by the projection to $P$.
If its image is trivial, multiply by $(1,g)^{-1}\in H$ to obtain
$(u,1)\in H$. This proves injectivity. Multiplication of lifts proves that
it is a group isomorphism. Its source is solvable, since
$\mathsf G_{u,\mathcal U}$ is unipotent, and its target is perfect, as a quotient
of $P$. A solvable perfect group is trivial: its derived series is
constant by perfectness and terminates by solvability.
Hence $A_{\rm u}=\mathsf G_{u,\mathcal U}$,
$A_P=P$, and the subdirect subgroup is the full product.
\end{proof}

\subsection{The logarithm and the ideal}
Put $K^D=K_{\rm full}T'$ and let $K_{\rm sol}$ be its solvable-coordinate
field over $T'$.
\begin{lemma}\label{lem:stabilizer}
For $u\in\Gal(\mathcal U/K_{\rm sol})$, put $u(L)=L+t(u)$. Then
\begin{equation}\label{eq:core-stabilizer}
\Gal(\Omega/K^D)=\{(u,\chi(u))\},\qquad
\chi(u)_{\mathfrak c}=\begin{cases}
\begin{pmatrix}1&-t(u)\\0&1\end{pmatrix},&\text{resonant},\\
1,&\text{nonresonant}.
\end{cases}
\end{equation}
\end{lemma}
\begin{proof}
Use the product coordinates of Proposition~\ref{prop:product}.
An automorphism fixing $K^D$ restricts to
$u\in\Gal(\mathcal U/K_{\rm sol})$.
Since $DL=1/z$, its action on $L$ is $L\mapsto L+t(u)$, with
$t(u)\in\C$. In a resonant representative, fixing $y$ and
$\Sigma=J+yL$ gives
\[
 y\longmapsto y,\qquad J\longmapsto J-t(u)y.
\]
In the ordered columns $(y,J)$ this is the matrix in
\eqref{eq:core-stabilizer}. In a nonresonant representative both
homogeneous solutions are fixed: the algebraic factor $z^{-2\nu}$
belongs to $T'$. Its matrix is therefore the identity.

Conversely, take any $u$ fixing $K_{\rm sol}$ and these class matrices.
The product formula supplies an automorphism of $\Omega$.
It fixes the representative coordinates by the preceding calculations.
For a composed resonant link,
$Qy_i=\kappa_1y_j$ and $QJ_i=\kappa_1(J_j+\theta y_j)$.
Here $\kappa_1\in T'$ and $\theta\in\C$ are fixed. Thus subtracting
$t(u)y_i$ is transported to subtracting $t(u)y_j$.
Thus the same shift fixes the entire parts in every linked block.
The nonresonant linkage coefficients belong to $T'$ and are fixed as
well. All derivative coordinates are fixed because the automorphisms
commute with $D$. These coordinates and $K_{\rm sol}$ generate $K^D$. This proves the converse and the stated graph.
\end{proof}
Projection onto $\Gal(\mathcal U/K_{\rm sol})$ is an isomorphism.
By (F5) and Lemma~\ref{lem:tower},
$\dim \mathsf G_{u,\mathcal U}-\dim\Gal(\mathcal U/K_{\rm sol})
=\trdeg_{T'}K_{\rm sol}=\rho+\tau$.
The finite algebraic extension $T'/T$ does not change this degree.
Applying (F5) once more gives
\begin{equation}\label{eq:core-degree}
\trdeg_{T'}K^D=\dim \mathsf G_{u,\mathcal U}+3\#\mathcal C-
\dim\Gal(\mathcal U/K_{\rm sol})=\rho+\tau+3\#\mathcal C.
\end{equation}
Algebraicity of $T'/T$ gives $\trdeg_TK_{\rm full}$, and the $d$
independent basis exponentials give the degree over $\C(z)$.
Over $T'$, the selected first-level basis generates the solvable
coordinates together with the higher levels. Each selected triple
recovers its fourth coordinate from the determinant equation.
The regular solution has value $1$ at zero. Thus division by it in the
function field is valid. The linkage maps recover the other blocks.
These recovery formulas show that the selected coordinates generate
$K^D$ over $T'$. Their number is $\rho+\tau+3\#\mathcal C$, which equals
its transcendence degree by \eqref{eq:core-degree}. A dependent generating
set would have smaller transcendence degree. They are therefore
independent over $T'$. Hence they are independent over $T$.

Adjoining $L$ adds the residue vector $e_{(0,0)}$.
The same calculation increases the degree by one exactly when
$e_{(0,0)}\notin\mathcal V_{\rm res}$. Then $L$ is transcendental over
$K_{\rm full}$, and the selected basis remains independent over $T(L)$.
Otherwise a constant combination of first-level coordinates has the
residue vector of $L$. Lemma~\ref{lem:primitive} expresses their
difference as an element of $R_0+\C$. Thus $L\in K_{\rm full}$.

The coefficient recurrences give
\begin{equation}\label{eq:core-inhomogeneous}
\mathscr B_{c,0}S_c=2I_c',\quad \mathscr K_{s,a}\Sigma_{s,a}=2F_{s,a}'-sF_{s,a}.
\end{equation}
For Bessel use $4(n+1)^2c_{n+1}=c^2c_n$ and $\mathsf H_{n+1}-\mathsf H_n=1/(n+1)$. For Kummer use $c_n=(a)_ns^n/(n!)^2$,
$(n+1)^2c_{n+1}=s(n+a)c_n$ and
$h_{n+1}-h_n=2/(n+1)-1/(n+a)$. Direct differentiation of $yL$
gives the same right-hand side as the corresponding normalized entire part.
Abel's formula and $y(0)=1$, $J\sim-L$ give Wronskians $-1/z$,
$-e^{sz}/z$. For nonresonant Bessel the series gives
$\operatorname{wr}(f_\nu,z^{-2\nu}f_{-\nu})=-2\nu z^{-2\nu-1}$.
These prove \eqref{eq:determinant-relations}. To obtain the common form,
keep the first three coordinates and use a distinct shifted fourth coordinate:
\[
 (X,Y,S,\widetilde W)=\begin{cases}
 (I,V,S_c,W_c-I/z),&\operatorname{Bess}(c,0),\\
 (F,F',\Sigma,\Sigma'-F/z),&\operatorname{Kum}(s,a),\\
 (f_\nu,f'_\nu,f_{-\nu},f'_{-\nu}-2\nu f_{-\nu}/z),
   &\operatorname{Bess}(c,\nu)\text{ nonresonant}.
 \end{cases}
\]
For Bessel, $\widetilde W=W_c-I/z$ and the raw coordinate is $W_c=S_c'$.
The other shifts are given in the display; the last subtracted term contains
$f_{-\nu}$. Each determinant becomes
\begin{equation}\label{eq:core-quadric}
X\widetilde W-YS+u/z=0,\qquad u=1,e^{sz},2\nu.
\end{equation}
For Bessel this is $I\widetilde W-VS_c+1/z=0$; the raw relation is
$z(IW_c-VS_c)-I^2+1=0$.

In the next lemma $W$ is an independent polynomial variable corresponding
to the shifted coordinate $\widetilde W$.
\begin{lemma}\label{lem:quadric}
Let $A$ be a domain and let $v$ be a unit of $A$.
The ring $Q_v=A[X,W,Y,S]/(XW-YS+v)$ is a free $A$-module whose
basis consists of the monomials not divisible by $XW$.
The natural map $A[X,Y,S]\to Q_v$ is injective.
If $A$ is a characteristic-zero field, $Q_v$ is geometrically integral
and has dimension three.
\end{lemma}
\begin{proof}
Use lexicographic order $X>Y>S>W$. The relation is monic with leading
monomial $XW$. Division by this monic polynomial gives unique normal forms over $A$. Thus the monomials not divisible by $XW$ form a free basis.
In particular the monomials in $X,Y,S$ remain independent.
Over an algebraic closure of a characteristic-zero field, a nontrivial
factorization would have two linear factors. Their quadratic part has
rank at most two. In contrast, $XW-YS$ has rank four.
The polynomial is therefore geometrically irreducible, and its quotient
is geometrically integral of dimension three.
\end{proof}
\begin{lemma}\label{lem:ideal}
The generators in Theorem~\ref{thm:main} generate the full kernels of
coordinate evaluation over $T$ and over
$\mathscr A_{\exp}=\C(z)[\Theta^{\pm1}]$. Both kernels are prime.
The same assertions hold over a coefficient field containing all data
of the generators and their fundamental residue constants.
\end{lemma}
\begin{proof}
A detailed proof is given in Appendix~\ref{app:detailed-proofs}.
All displayed generators vanish on the functions. First eliminate duplicate
and constant slots using their monic graphs. The exponential, residue and linkage generators are also graphs. Each eliminates one variable in terms of retained variables. Its coefficient on that variable is $1$.
Successive substitution therefore gives isomorphisms of quotient rings.
After these eliminations and the invertible fourth-coordinate shifts
in \eqref{eq:core-quadric}, the quotient is
\begin{equation}\label{eq:core-relation-ring}
 T[Z_1,\ldots,Z_{\rho+\tau}]
 \otimes_T\bigotimes_{\mathfrak c}Q_{u_{\mathfrak c}/z}.
\end{equation}
Here the $Z_i$ are the residue-basis and distinct higher-tower variables.
Lemma~\ref{lem:quadric} and (F13) make this a domain of dimension
$\rho+\tau+3\#\mathcal C$.
Evaluation maps this quotient to the coordinate algebra, whose fraction
field is $K_{\rm full}$. Its kernel is prime, and the image has the same
dimension by \eqref{eq:core-degree}. If the kernel were nonzero, its prime would strictly contain the zero prime. By (F13), the quotient dimension would decrease. Thus evaluation is injective. This proves the
ideal assertion over $T$.

Over $\mathscr A_{\exp}=\C(z)[\Theta^{\pm1}]$, the identical graph eliminations are
$\mathscr A_{\exp}$-linear, and every $u_{\mathfrak c}/z$ remains a unit. The quotient
has the same tensor presentation over $\mathscr A_{\exp}$. The polynomial factor has
its monomial basis, and each quadric factor has the free basis of
Lemma~\ref{lem:quadric}. Products of these basis elements form an
$\mathscr A_{\exp}$-basis of the tensor product.
Write $I$ for the generated ideal and
$B=\mathscr A_{\exp}[X]/I$ for this quotient. If a class $b\in B$ becomes
zero in its generic fibre, some nonzero $a\in\mathscr A_{\exp}$ satisfies
$ab=0$. Freeness forces $b=0$. Thus $B$ injects into that fibre.
By Lemma~\ref{lem:exponential-log}, $\Theta_h\mapsto t_h$ identifies
$\operatorname{Frac}\mathscr A_{\exp}$ with $T$. The generic fibre is
exactly \eqref{eq:core-relation-ring}. Evaluation on it is injective by
the preceding paragraph. Any Laurent polynomial vanishing on the functions has zero class in the generic fibre. Hence its class in $B$ is zero. It belongs to $I$. This proves that the Laurent evaluation kernel is
exactly the stated ideal.
The quotient embeds in a field. Thus it is a domain. Its fraction field
has degree $d+\rho+\tau+3\#\mathcal C$ over $\C(z)$, giving the
asserted prime dimension.

For clarity, its contraction is an explicit elimination ideal. Introduce
new variables $V_h$ and equations $\Theta_hV_h-1$. Rewrite all Laurent
generators using $V_h$ for $\Theta_h^{-1}$, obtaining an ideal $J$ in
$\C(z)[X,\Theta,V]$. Then
\[
 \ker\bigl(\C(z)[X]\longrightarrow K_{\rm coord}\bigr)
   =J\cap\C(z)[X].
\]
This equality follows from the quotient isomorphism with $\mathscr A_{\exp}[X]$ and
restriction of evaluation. Equivalently, clear Laurent denominators,
saturate by $\prod_h\Theta_h$, and eliminate $\Theta$. Substituting listed
exponential variables requires saturation by their product as well.
These inverses encode the fact that every exponential is nonzero.

Let $k_1$ contain all coefficients, listed constants and fundamental residue constants.
The same free bases define the quotient over
$\mathscr A_{\exp,k_1}=k_1(z)[\Theta^{\pm1}]$.
The graph eliminations and monic divisions use only these coefficients.
Thus extension to $\mathscr A_{\exp}$ sends each basis element to the identical basis
element. As $\mathscr A_{\exp,k_1}\hookrightarrow \mathscr A_{\exp}$, this extension is injective.
Composing with evaluation over $\mathscr A_{\exp}$ embeds the original quotient in the
function field. Its defining ideal is therefore the contraction of the
complex-coefficient kernel and is prime.

The free variables and every triple $(X,Y,S)$ remain independent over
$\mathscr A_{\exp,k_1}$ by the same monomial basis. In the fraction field the fourth
coordinate is $W=(YS-u/z)/X$. The fraction field is consequently
rational in these selected variables over $k_1(z)(\Theta)$.
After evaluation it is $k_1(z)(t,\mathcal F)$, with degree
$d+\rho+\tau+3\#\mathcal C$. This proves coefficient-field descent
and the stated degree.

For compatible input, Appendix~\ref{app:algorithm} points to the two
formulas of Supplementary Section~\ref{sup:input-model}. They put the fundamental residue constants in
$k_1$. Each chosen basis relation uses one normalization. Arbitrary
relations are combinations of them. Each listed constant adds $X_c-c$.
Duplicate graphs are retained before closure and in every saturation or
contraction. Eliminating an auxiliary representative transfers its graph
to the retained original slots by substitution.
\end{proof}

\section{The linkage classification}\label{sec:linkage}

Gorelov classified cogredience and contragredience, equivalently rank-one-twist
isomorphism, of irreducible equations $L(\nu;\lambda;\alpha z^p)y=0$
\cite[Lemma~12 and proof, p.~7; Lemma~14, p.~7, proof pp.~8--9]{Gorelov2021}.
The mixed Bessel--Kummer case is condition~(9)
\cite[p.~3]{Gorelov2021}. Its sufficiency is
\cite[Theorem~2, condition~(5), p.~138; proof pp.~141--142]{Gorelov2020}.
Theorem~\ref{thm:linkage} specializes this classification.
Below we prove normalized E-part maps and composition for the classical regular-solution transformations.

For these links, the regular solutions and their normalized E-parts satisfy
\begin{equation}\label{eq:link-kummer}
\begin{aligned}
 f_a(u)&=e^u f_{1-a}(-u),\\
 S_a(u)&=e^u S_{1-a}(-u),
\end{aligned}
\end{equation}
\begin{equation}\label{eq:link-contiguity}
\begin{aligned}
 a f_{a+1}(u)&=u f'_a(u)+af_a(u),\\
 aS_{a+1}(u)&=uS'_a(u)+aS_a(u)-f_a(u)+f_{a+1}(u),
\end{aligned}
\end{equation}
\begin{equation}\label{eq:link-half}
\begin{aligned}
 f_{1/2}(2cz)&=e^{cz}I_c(z),\\
 S_{1/2}(2cz)&=e^{cz}S_c(z).
\end{aligned}
\end{equation}
For $\nu\in\Q\setminus\Z_{\le0}$, the series-defined functions satisfy
\begin{equation}\label{eq:link-order}
 f_{\nu+1,c}=\frac{2(\nu+1)}{c^2z}f'_{\nu,c},\qquad
 f_{\nu-1,c}=f_{\nu,c}+\frac{z}{2\nu}f'_{\nu,c},
\end{equation}
and integer orders are given by Lemma~\ref{lem:integer}. Write
$\operatorname{Sol}(M)$ for the solution space of $M$. For an oriented
resonant link, let $\Phi_{ij}:\operatorname{Sol}(M_i)\to
\kappa_{ij}\operatorname{Sol}(M_j)$ be the induced map, with
$\kappa_{ij}=\kappa_1$ and the convention
\begin{equation}\label{eq:link-composition-convention}
\Phi_{ik}=\kappa_{ij}\Phi_{jk}\kappa_{ij}^{-1}\circ\Phi_{ij}.
\end{equation}
Multiplication by $\kappa_{ij}$ identifies the target with the twisted
solution space. When $\Gamma$ is saturated for the list,
$\kappa_{ij}\in T^\times$. For every composed resonant link, written as $Q$,
\begin{equation}\label{eq:link-composed}
 Qy_i=\kappa_1y_j,\qquad QJ_i=\kappa_1(J_j+\theta y_j),\quad\theta\in\C.
\end{equation}
Here $(\kappa_1,\theta)$ is $(e^{sz},0)$, $(1,-1/a)$, or
$(e^{cz},0)$ for the following oriented maps. The first is the identity from
$\operatorname{Kum}(s,a)$ to $\operatorname{Kum}(-s,1-a)$.
The second is the contiguity operator $Q=(zD+a)/a$ from $\operatorname{Kum}(s,a)$ to
$\operatorname{Kum}(s,a+1)$. The third is the identity from
$\operatorname{Kum}(2c,1/2)$ to $\operatorname{Bess}(c,0)$.
For the identical order-zero Bessel operators, $(\kappa_1,\theta)=(1,0)$.
Along a path the twist factors multiply and the constants add. Reversing
an elementary link replaces them by $(\kappa_1^{-1},-\theta)$.
For a nonresonant link, the rational gauges and the algebraic power twist
$z^\upsilon\in T'=T(z^{1/\ell})$ give the corresponding map over $T'$. The twist need not belong to $T$. The listed $f_{-\nu,c}$ omits the
factor $z^{-2\nu}$. Thus induced coordinate maps lie over $k\subset T$.
There is no logarithmic solution $J$
in this case, and $\ell$ clears the denominators of $2\nu_0$.
Differentiation gives the derivative-coordinate maps, using Section~\ref{sec:setup}'s convention.
\begin{proof}[Proof of Theorem~\ref{thm:linkage}]
The parameter substitution in Appendix~\ref{app:gorelov} identifies the
homogeneous cases with Gorelov's classification. The same-power case of
\cite[Lemma~14, using Lemma~12, p.~7]{Gorelov2021} gives
\eqref{eq:link-bb}--\eqref{eq:link-kk}. Its mixed condition~(9)
\cite[p.~3]{Gorelov2021} gives
$2\nu_0\in\Z$, $s^2=4c^2$ and $2a\in\Z$. Thus
\eqref{eq:link-bk} follows. The mixed sufficiency is also
\cite[Theorem~2, p.~138]{Gorelov2020}. Lemmas~\ref{lem:formal}
and~\ref{lem:twist} give the rates and powers.

For the contiguity calculation assume $a\notin\Z$ and take $s=1$. Thus $u=z$ and $D=d/du$.
The operator $(uD+a)/a$ has matrix and inverse
\begin{equation}\label{eq:contiguity-matrix}
 \frac1a\begin{pmatrix}a&u\\a&u+a\end{pmatrix},\qquad
 \frac1a\begin{pmatrix}u+a&-u\\-a&a\end{pmatrix}.
\end{equation}
Its determinant is $1$. A nonzero kernel element would be proportional to $u^{-a}$.
This function is no solution for $a\notin\Z$. For the other links, use
\[
 \mathscr K_{s,a}(e^{sz}y)=e^{sz}\mathscr K_{-s,1-a}y,\qquad
 \mathscr K_{2c,1/2}(e^{cz}y)=e^{cz}\mathscr B_{c,0}y.
\]
Normalization at zero gives the regular Kummer identities.
For $f_{\pm1/2,c}$ and the order maps, compare the defining series directly:
$n!(3/2)_n=(2n+1)!/4^n$ and $n!(1/2)_n=(2n)!/4^n$ give
$f_{1/2,c}=\sinh(cz)/(cz)$ and $f_{-1/2,c}=\cosh(cz)$;
normalization at zero alone does not fix the order $-1/2$ solution.
Changing $\nu$ to $-\nu$ introduces $z^{2\nu}$. The recurrences
\eqref{eq:link-order} and Lemma~\ref{lem:integer} give the order maps.

To obtain the normalized E-parts, write $J=\Sigma-yL$.
Monodromy preserves the invariant line. Thus the image of $J$ has the form
\eqref{eq:link-composed}. Under the two exponential conjugations, the
E-parts have equal forcing. Their difference is a homogeneous holomorphic solution vanishing at zero. This proves the corresponding identities. For contiguity, direct differentiation gives
\[
 \frac{uD+a}{a}(S_a-f_aL)
 =\frac{uS'_a+aS_a-f_a}{a}-f_{a+1}L,
\]
whose holomorphic part at zero is $-1/a$. Equivalently,
$a(a+1)_n=(a+n)(a)_n$ and
$h_n(a+1)=h_n(a)+1/a-1/(a+n)$. These prove the E-part maps and
constants. For composition, suppose $Q^{(1)}$ gives
$\kappa_1(y_m,J_m+\theta_1y_m)$ and $Q^{(2)}$ gives
$\kappa_2(y_j,J_j+\theta_2y_j)$. The transported operator in
\eqref{eq:link-composition-convention} sends the pair to
$\kappa_1\kappa_2(y_j,J_j+(\theta_1+\theta_2)y_j)$.
The inverse transported map is
$Q^{(1),\rm rev}(v)=\kappa_1^{-1}(Q^{(1)})^{-1}(\kappa_1v)$, which gives
$(\kappa_1^{-1},-\theta_1)$. This proves both orientation rules.
\end{proof}

The regular identities are DLMF~13.3.17, 13.2.39 and 10.39.5.
After normalization, DLMF~13.2.9, 13.2.41 and 13.6.10 also give the E-parts \cite{DLMF}.

\section{Values at algebraic points}\label{sec:values}
Use the arithmetic hypotheses and notation of Theorem~\ref{thm:values}.
The $(\mathrm{B2}^E)$ residue vectors have
$\sum_\mu v_{(\mu,0)}=0$, which implies $e_{(0,0)}\notin\mathcal V_{\rm res}$.
Lemma~\ref{lem:integer} reduces integer-order blocks; Supplementary Section~\ref{sup:classical-comparison} compares earlier value inputs.

\begin{proof}[Proof of Theorem~\ref{thm:values} and Corollary~\ref{cor:value-base}]
A detailed proof is given in Appendix~\ref{app:detailed-proofs}.
The finite list's data lie in one number field.
Appendix~\ref{app:arithmetic}'s estimates give the $E$-function arithmetic bounds
after order-zero reduction. The primitives' entire derivatives and
zero-origin normalization remove the apparent pole at zero.

Consider the $\Qbar(z)$-linear system with coordinates
\begin{equation}\label{eq:value-system}
\begin{gathered}
1,\quad e^{\pm\lambda_hz},\quad e^{\mu z}\ (\mu\text{ kernel rate or }-\eta),\\
\text{all primitives, tower levels and block coordinates}.
\end{gathered}
\end{equation}
Its only possible finite pole is zero. In addition to the recurrences
already given, its equations include
\[
I''=c^2I-I'/z,\quad S''=c^2S-S'/z+2I'/z,\quad
F''=(s-z^{-1})F'+asF/z,
\]
\[
\Sigma''=(s-z^{-1})\Sigma'+as\Sigma/z+(2F'-sF)/z,\quad
f''_{\nu,c}=c^2f_{\nu,c}-(2\nu+1)f'_{\nu,c}/z.
\]
Successive-derivative dependence gives scalar equations over $\Qbar[z]$,
completing the $E$-function hypotheses of (F1).
Every $\alpha\ne0$ is ordinary. Put $R=\Qbar[z^{\pm1}][\Theta^{\pm1}]$.
By \eqref{eq:link-bb}--\eqref{eq:link-bk}, two blocks of one class are
joined by a finite path of the displayed maps.
Scale sign $c\mapsto-c$ and order sign $\nu\mapsto-\nu$ use
identities or permutations.
Order shifts use \eqref{eq:link-order} and Lemma~\ref{lem:integer}.
Contiguity and its inverse have determinant $1$.
The Kummer transformation and half link multiply by $e^{\pm sz}$ and
$e^{\pm cz}$. These are monomials in $\Theta$ because the rates lie
in $\Gamma$.
The corresponding E-part identities are
\eqref{eq:link-kummer}--\eqref{eq:link-half}.
Each map and its inverse on $(y,y',\Sigma,\Sigma')$ has entries in $R$
and determinant a unit of $R$. Hence every type-4 generator lies in $R[X]$.
For the shift $\nu\to\nu+1$, the $2\times2$ map on $(f_{\nu,c},f'_{\nu,c})$ has
determinant $-4(\nu+1)^2/(c^2z^2)$, whereas the full $4\times4$
coordinate map has determinant $(\nu+1)^2/\nu^2$.
Normalization excludes zero parameter factors, so each inverse specializes at $\alpha$.
Linkage elimination inverts no coordinate value.

For type 2, $\kappa_{ji}\in\Qbar$ since residue vectors are algebraic;
\[
 D\Bigl(F_j-\sum_i\kappa_{ji}F_i\Bigr)
   =\sum_\mu q_\mu e^{\mu z},\qquad
 q_\mu\in\Qbar[z^{\pm1}],\qquad
 \Res_0 e^{\mu z}q_\mu=0.
\]
On $\Qbar[z^{\pm1}]$, the map $g\mapsto g'+\mu g$ has image exactly
the kernel of $q\mapsto\Res_0e^{\mu z}q$ (triangular on $z^n$, $n\ge0$,
and $z^{-m-1}\equiv(\mu/m)z^{-m}$ modulo the image), so $r_j\in R$.
All first-level coordinates vanish at zero, hence
$c_j=-\lim_0r_j\in\Qbar$.
This finite Laurent--exponential constant coefficient keeps the affine graph in $R[X]$.
Types 1, 3, half-order, duplicate and constant graphs visibly lie in $R[X]$.

Duplicate and constant graphs specialize unchanged. Integer-order blocks
are first replaced by the order-zero blocks in Lemma~\ref{lem:integer}. Its matrix and inverse are regular at $\alpha\ne0$. Adjoin these blocks,
apply the following argument, and eliminate them afterwards.
After graph elimination at $\alpha$, the determinant shifts are regular.
They give
\begin{equation}\label{eq:specialized-quadric}
\begin{gathered}
R_\Theta[X]/\mathfrak P_\alpha\simeq
R_\Theta[\rho^E+\tau\text{ free variables}]\otimes_{R_\Theta}
\bigotimes_{\mathfrak c}Q_{u_{\mathfrak c}/\alpha},\\
u_{\mathfrak c}/\alpha\in\{1/\alpha,\Theta^{v_s}/\alpha,2\nu/\alpha\}.
\end{gathered}
\end{equation}
Lemma~\ref{lem:quadric} gives freeness over $R_\Theta$ and injects the selected polynomial subring.
Freeness prevents localization torsion, so the quotient injects into its generic fibre.
The quadric parameters are units; Lemma~\ref{lem:quadric} and (F13) make
that fibre geometrically integral. The quotient is therefore a domain of
dimension $d+\rho^E+\tau+3\#\mathcal C$. Its generators vanish under evaluation.
Theorem~\ref{thm:main} and (F1) on \eqref{eq:value-system} give the same dimension for $I_\alpha$.
Extra system exponentials are monomials in the basis exponentials.
Thus $\mathfrak P_\alpha\subset I_\alpha$ are prime with equal quotient
dimensions; (F13) excludes strict inclusion.
The selected subring still injects under evaluation, proving independence
and \eqref{eq:value-nonzero} without division by the regular coordinate.

A relation over $\mathcal E$ has coefficients in a finite algebraic
extension of $\Qbar(e^{\beta_1},\ldots,e^{\beta_s})$. Adjoin the rates
$\beta_i/\alpha$ to $\Gamma$, saturate for the enlarged list, and list
the resulting basis exponentials.
Its residue rank, higher-level count and classes are unchanged. The proved
result, localized over $K_E=\Qbar(e^{\lambda'_h\alpha}:h)$ for the
enlarged rate basis $(\lambda'_h)$, gives the same degree.
It contains every $e^{\beta_i}$; only the exponential basis and count change.
In $K_E[X]$, substitution of each exponential variable by its value is a
retraction to $K_E[X_{\rm ne}]$. It kills the type 1 generators. Thus the
contracted ideal is generated by types 2--4, the half-integer exponential
graphs and their derivative and order-shift graphs, duplicate graphs and
listed-constant graphs. These are the stated nonexponential generators. The generic fibre is
geometrically integral. Let $K'/K_E$ be a finite algebraic extension
containing the relation's coefficients. The extended ideal in
$K'[X_{\rm ne}]$ remains prime by geometric integrality. Equivalently,
it is the contraction after specializing the exponential variables in
$K'[X]$. Its evaluation kernel in $K'[X_{\rm ne}]$ has the same dimension;
(F13) proves equality: algebraic coefficients change neither transcendence
degree nor the geometrically integral model's dimension.
Every relation over $\mathcal E$ occurs in such a finite enlargement.
This proves \eqref{eq:value-exponential-base}
and its ideal assertion. For an infinite family, a polynomial relation involves finitely many values. Close that finite list by adding lower tower levels and the
four coordinates of each chosen block. The listed first-level residue
vectors are independent. Thus extend them to a residue basis of the closure.
The higher-level pairs remain distinct, and the chosen triples represent
different classes. Thus the original finite family is a subset of the
selected independent set for that closure. The finite theorem excludes
the proposed relation. For values $q_i(\alpha_i)$ at several $\alpha_i\in\Qbar^\times$,
replace $q_i$ by $q_i(\alpha_i z)$. Nodes, scales and rates are then multiplied
by $\alpha_i$, and $(\mathrm{B2}^E)$ kernels $q(z)$ become
$\alpha_iq(\alpha_i z)$. A residue indexed by $(\mu,p)$ is
reindexed by $(\alpha_i\mu,p/\alpha_i)$ with the same value, and each derivative
coordinate is multiplied by $\alpha_i$. Apply the selection conditions
(residue independence, distinct $(\eta\alpha_i,k)$, one triple per class)
to the rescaled list at $z=1$. The derivative factors are invertible
algebraic scalars. Joint closure and saturation retain any collisions
introduced by rescaling before the theorem is applied.
\end{proof}

\begin{lemma}\label{lem:integer}
For $n\ge1$,
\begin{equation}\label{eq:integer-bessel-matrix}
\begin{gathered}
\begin{pmatrix}I_{n,c}\\I'_{n,c}\\P_{n,c}\\P'_{n,c}\end{pmatrix}
=\begin{pmatrix}M_n&0\\N_n&M_n\end{pmatrix}
\begin{pmatrix}I_c\\V_c\\S_c\\W_c\end{pmatrix},\\
M_n,N_n\in M_2(\Q[c^{\pm1},z^{-1}]),\qquad\det M_n=(-1)^n.
\end{gathered}
\end{equation}
It is invertible at every nonzero algebraic $\alpha$ when $c$ is algebraic.
\end{lemma}
\begin{proof}
From \eqref{eq:bessel-integer},
\[
M_n=\begin{pmatrix}a_n&b_n\\a'_n+c^2b_n&a_n+b'_n-b_n/z\end{pmatrix}.
\]
Differentiating $P_{n,c}$ with $S''=c^2S-W/z+2V/z$ gives $N_n$.
The recurrence \mbox{sends} $(I_0,K_0)$ to $(I_n,(-1)^nK_n)$. Their Wronskians
are $-1/z$ and $(-1)^n(-1/z)$ \cite[10.28.2,10.29.1(a),10.29.3(a)--(b)]{DLMF}. This proves the
determinant. For example,
\[ M_1=\begin{pmatrix}0&1/c\\c&-1/(cz)\end{pmatrix}. \]
The integer-order and order-zero blocks therefore generate the same
fields. To obtain the integer-order ideal, adjoin the order-zero block,
specialize this invertible map, and eliminate the added coordinates.
This procedure also applies when $P_{n,c}$ has a pole at zero.
\end{proof}

\section{Mixed connection values}\label{sec:gamma}
All nodes in this section are positive algebraic numbers; all logarithms
are real. Put $\mathcal E=\operatorname{acl}_{\C}\Qbar(e^\beta:\beta\in\Qbar)$.
Write $R_B(u)=\mathsf S(u)/I_0(u)$ and $R_a(t)=S_a(t)/f_a(t)$ for the
argument-normalized functions of Section~\ref{sec:setup}.
For $a\in\Q\setminus\Z$, write $a=\vartheta_a+k_a$, where
$\vartheta_a=p_a/q_a\in(0,1)$ is reduced and $k_a\in\Z$. Define
\[
 G_\vartheta=\psi(\vartheta)+\gamma,\qquad
 H(\vartheta,k)=\begin{cases}
 \sum_{j=0}^{k-1}(\vartheta+j)^{-1},&k\ge0,\\
 -\sum_{j=k}^{-1}(\vartheta+j)^{-1},&k<0.
 \end{cases}
\]
\begin{theorem}[Mixed connection values]\label{thm:connection}\label{ct2:statement-start}
Select $N\ge1$ rows of the following table, with positive algebraic
nodes $x,u,t$, integer orders $m\ge0$ and rational nonintegral parameters $a$.
Require distinct Ein nodes and distinct Bessel nodes.
Exclude Kummer pairs at the same $t$ whose parameters differ by an integer.
Also exclude Bessel--Kummer pairs with $t=2u$ and $a\in\tfrac12+\Z$.
Put $\epsilon_m=(-1)^m$ and define
\[
\begin{gathered}
 (p_0,v_0)=(1,0),\quad(p_1,v_1)=(0,1),\\
 (p_{m+1},v_{m+1})=(p_{m-1},v_{m-1})-\frac{2m}{u}(p_m,v_m),
 \quad\beta_m=-\frac{\epsilon_mv_m}{uI_0(u)I_m(u)}.
\end{gathered}
\]
The denominators below are nonzero. The table defines $V=X-d$:
\begin{equation}\label{eq:ct2-table}
\begin{array}{c|c|c}
 V_i&X_i&d_i\\ \hline
 E_1(x)&\Ein(x)&\gamma+\log x\\
 K_m(u)/I_m(u)&\epsilon_mR_B(u)+\beta_m(u)&
       \epsilon_m\{\gamma+\log(u/2)\}\\
 \Gamma(a)U(a,1,t)/f_a(t)&R_a(t)-H(\vartheta_a,k_a)&
       \gamma+\log t+G_{\vartheta_a}
\end{array}
\end{equation}
For an arbitrary integer Bessel order, use its absolute value. Set
\[
 B=\operatorname{acl}_{\C}\mathcal E
 \bigl(I_0(u),I_1(u),f_a(t),f'_a(t):\text{selected rows}\bigr).
\]
Gauss's formula and multiplicative relations among the positive algebraic
arguments of the logarithms express $d=A\mathbf c$, with $A\in M_{N,r_{\mathbf c}}(\Qbar)$
and $r_{\mathbf c}$ the number of retained constants.
Here $\mathbf c$ consists of $\gamma$, the required real logarithms, and $\pi$
only if its coefficient is nonzero; logarithms are combined and unused entries omitted.
For each $F=\mathcal E,B$,
\begin{equation}\label{eq:ct2-actual-degree}
 \trdeg_F F(\mathbf c,V_1,\ldots,V_N)\ge N.
\end{equation}
Thus these numbers contain an algebraically independent subset of
size $N$ over either base field. If all entries of $\mathbf c$ are algebraic over
$F$, the values $V_1,\ldots,V_N$ themselves are algebraically independent
 over $F$.
\label{ct2:statement-end}
\end{theorem}
\begin{proof}\label{ct2:proof-start}
The connection expansions \cite[6.2.4, 10.28.2, 10.31.1, 13.2.9]{DLMF} give
\begin{align*}
 E_1(x)&=\Ein(x)-\gamma-\log x,\\
 K_m(u)/I_m(u)&=\epsilon_m\{R_B(u)-\gamma-\log(u/2)\}+\beta_m(u),\\
 \Gamma(a)U(a,1,t)&=S_a(t)-f_a(t)\{\log t+\psi(a)+2\gamma\}.
\end{align*}
The Bessel identity follows from $I_m=p_mI_0+v_mI_1$,
$\epsilon_mK_m=p_mK_0-v_mK_1$ and $I_0K_1+I_1K_0=1/u$.
In particular $I_m(u)>0$ and $\beta_m(u)\in B$.
Theorem~\ref{thm:values}, applied to the closed normalized Kummer block
at $z=1$, gives $f_a(t)\ne0$, including for negative rational $a$.
This single-block class is its own representative, so its regular
coordinate is $f_a(t)$ in \eqref{eq:value-nonzero}.
The recurrence gives $\psi(a)+\gamma=G_{\vartheta_a}+H(\vartheta_a,k_a)$.
Gauss's formula is
\begin{equation}\label{eq:ct2-gauss}
 G_{p/q}=-\log(2q)-\frac\pi2\cot\frac{\pi p}{q}
 +2\sum_{j=1}^{\lfloor(q-1)/2\rfloor}
       \cos\frac{2\pi pj}{q}\log\sin\frac{\pi j}{q}.
\end{equation}
Its coefficients and positive logarithm arguments are algebraic.
Taking a multiplicative $\Q$-basis of the arguments proves $d=A\mathbf c$
with the stated reduced constant column.

The exclusions place the selected second-order blocks in distinct linkage classes.
Distinct Ein nodes have independent residues $e_{(0,0)}-e_{(-x,0)}$.
Apply closure, normalization and rate saturation.
By Theorem~\ref{thm:values} and Corollary~\ref{cor:value-base}, the Ein values and one triple
$(I_0,I_1,\mathsf S)$ or $(f_a,f'_a,S_a)$ per selected class are then jointly independent
 over $\mathcal E$. Adjoin the regular pairs and divide by their nonzero
regular coordinates. Then pass to the algebraic closure $B$.
The remaining $N$ ratios are independent. Nonzero diagonal scalings and
$B$-valued shifts show that $X_1,\ldots,X_N$ are independent over $B$,
hence over $\mathcal E$.
Since $X=V+A\mathbf c$ and $A$ is algebraic, $F(X)\subset F(\mathbf c,V)$ proves
\eqref{eq:ct2-actual-degree}. A finite generating set contains a
transcendence basis. If $\mathbf c$ is algebraic over $F$, it lies in $F$,
since both permitted base fields are algebraically closed within $\C$.
Then $V=X-A\mathbf c$ is an $F$-valued translation of $X$.
\label{ct2:proof-end}
\end{proof}

\begin{corollary}\label{cor:gamma-gompertz}
At least one of $\gamma$ and $E_1(1)$, equivalently at least one of
$\gamma$ and $\delta=eE_1(1)=\int_0^\infty e^{-s}/(1+s)\,ds$,
is transcendental over $\mathcal E$.
\end{corollary}
\begin{proof}
Use the Ein row at $x=1$: $B=\mathcal E$, $\mathbf c=(\gamma)$ and
$\Ein(1)=\gamma+E_1(1)$. Also $e\in\mathcal E^\times$.
\end{proof}
The stated alternative over the full field $\mathcal E$ is known.
It follows from Shidlovskii's Chapter~7, \S1, Theorem~1
\cite[pp.~232--233]{Shidlovskii1989}, applied after the finite-rate
reduction in Supplementary Section~\ref{sup:classical-comparison}.
Mahler announced an exponential-integral statement without proof
\cite[\S23]{Mahler1968}. Rivoal proves the one-point result for
algebraic $x\notin(-\infty,0]$ \cite[Theorem~2(ii), \S7]{Rivoal2012}.
Theorem~\ref{thm:connection} assembles multi-node mixed families over $B$.
The comparison with the sources is partial.

\begin{example}\label{ex:ct2-mixed}
Choose $x=1$, $(u,m)=(2,0)$ and $(t,a)=(1,1/2)$.
The rows are admissible since the mixed linkage would require $t=4$.
Here $\mathbf c=(\gamma,\log2)^t$ and
\[
 B_*=\operatorname{acl}_{\C}\mathcal E
 (I_0(2),I_1(2),f_{1/2}(1),f'_{1/2}(1)).
\]
For each $F=\mathcal E,B_*$, the five numbers
\[
 \gamma,\quad\log2,\quad E_1(1),\quad
 \frac{K_0(2)}{I_0(2)},\quad
 \frac{\Gamma(1/2)U(1/2,1,1)}{f_{1/2}(1)}
\]
contain three algebraically independent numbers over $F$.
Supplementary Section~\ref{sup:cancellation-remarks} computes $A$ and an explicit cancellation.
\end{example}

\begin{proposition}[Coefficient cancellation]\label{prop:coefficient-cancellation}\label{lem:coefficient-cancellation}
Let $\kappa$ be the dimension of $\Span_B\{d_i+B\}$ in $\C/B$.
Let $J,J_0$ have row bases of $\{\lambda:\lambda d\in B\}$ and
$\{\lambda:\lambda d=0\}$ in $B^{1\times N}$.
For independent indeterminates $Z$, put
$\mathcal A_d=\{P\in B[Z]:P(Z-d)\in B[Z]\}$. Then
\begin{equation}\label{eq:ct2-algebra}
\begin{gathered}
 \mathcal A_d=B[JZ],\qquad
 \{P\in B[Z]:P(Z-d)=P(Z)\}=B[J_0Z],\\
 B[JV]=B[JX],\qquad\trdeg_B B(JV)=N-\kappa,\qquad
 N-\rank_{\Qbar}A\le N-\kappa\le N.
\end{gathered}
\end{equation}
\end{proposition}
Supplementary Section~\ref{sup:cancellation-remarks} proves the proposition in full.
Neither $\kappa$ nor $J$ is claimed effective. Even whether $\gamma+\log x\in B$ is undecided.
Accidental intersections $B[V]\cap B[X]$ and equalities $P(V)=P(X)$ are not classified.
No independence of $\mathbf c$ or of $X$ over $B(\mathbf c)$ is assumed.
Linked rows must be reduced first; higher Ein tails are excluded.
The irrationality of Euler's constant remains open.
\label{ct2:section-end}


\section*{Use of AI tools}
The author devised the methods and the main arguments. Codex (OpenAI, GPT-6 Astra)
checked the derivations and the cited sources, and ran the exact computations of the
supplementary sections. Claude Code (Anthropic, Claude Fable 5) checked the arguments
and the references. The author checked every step of the proofs, every computation and
every reference, and takes full responsibility for the content of the paper.

\section*{Declarations}

\textbf{Competing interests.} The author has no competing interests.

\textbf{Funding.} No funding was received for this work.

\textbf{Data availability.} The ancillary script and its exact output are in
\texttt{anc/} within the source directory; see \texttt{anc/N1-README.md}.

\label{page:body-end}
\clearpage
\appendix
\label{page:appendix-start}
\section{Foundational facts}\label{app:foundations}\label{subsec:foundations}
An $E$-function is an entire series $\sum a_nz^n/n!$ with algebraic
coefficients, $h(a_0,\ldots,a_n)=O(n)$ and a linear differential equation
over $\Qbar[z]$ \cite[p.~369]{Beukers2006}.
Here $h$ is the absolute logarithmic Weil height of $[1:a_0:\cdots:a_n]$.
Our differential fields have characteristic zero and common algebraically
closed constants $C_0$.
\newcounter{foundation}\renewcommand{\thefoundation}{F\arabic{foundation}}
\newcommand{\foundation}[1]{\par\noindent\refstepcounter{foundation}(\thefoundation)\label{fact:#1}\enspace}
\foundation{F1}For $E$-functions $f'=Af$, $A\in M_n(\Qbar(z))$, and a
common denominator $Q$ of $A$, every algebraic $\alpha$ with
$\alpha Q(\alpha)\ne0$ satisfies
\[
\trdeg_{\Qbar}\Qbar(f(\alpha))=\trdeg_{\Qbar(z)}\Qbar(z)(f)
\]
\cite[Theorem~1.1, p.~370]{Beukers2006}.
\foundation{F2}Under (F1), each homogeneous relation $P(f(\alpha))=0$
lifts to $R\in\Qbar[z,X]$ homogeneous in $X$, with
$R(z,f)=0$, $R(\alpha,X)=P$. Adjoin $1$ for inhomogeneous relations
\cite[Theorems~1.2--1.3 and the intervening paragraph, p.~370]{Beukers2006}.
\foundation{F3}If primitives with $y_i'\in F_{\rm b}$ are algebraically dependent
over $F_{\rm b}$, a nonzero constant combination of them belongs to $F_{\rm b}$
\cite[pp.~315--316]{Ostrowski1946}. For the differential-algebraic setting, see \cite{Kolchin1973};
Lemma~\ref{lem:primitive} proves the same-constant meromorphic case.
\foundation{F4}A subdirect Lie algebra in $\bigoplus\mathfrak{sl}_2$ is a
sum of diagonal graphs (Lemma~\ref{lem:goursat} and
Proposition~\ref{prop:product}). The factor automorphisms are inner,
as proved below.
\foundation{F5}A fundamental matrix generates a PV extension with
$\trdeg(E/K)=\allowbreak\dim\Gal(E/K)$. Closed subgroups correspond to
intermediate fields, normal subgroups to PV subextensions, and invariant
subspaces for the full group to differential submodules. Our lower-unitriangular
matrices satisfy $\sigma(F)=FC_\sigma$, so $C_\sigma=F^{-1}\sigma(F)$
is lower unitriangular. The PV group is algebraic
\cite[Observation~1.26(2), p.~19; Theorem~1.27(1), p.~20]{VanDerPutSinger2002Draft}.
Each nonidentity element's powers have connected additive closure
\cite[\S7.3, Remark, p.~106]{Borel1991}, proving unipotence and connectedness.
Liouvillian PV
extensions are exactly those with solvable Galois identity component \cite[Theorem~1.43, pp.~33--34]{VanDerPutSinger2002Draft}.
\foundation{F6}The equation $x^2y''+xy'+(x^2-\nu^2)y=0$ has a nonzero
Liouvillian solution exactly for $\nu\in\tfrac12+\Z$
\cite[\S3.2, Example~2, p.~14; \S4.2, Example~2, pp.~19--20]{Kovacic1986}.
Nonzero scaling and the algebraic gauges used here preserve this property
by direct change of variables.
\foundation{F7}The Laurent, logarithmic-derivative and polynomial-log
reductions used here are proved in Lemmas~\ref{lem:reduction},
\ref{lem:exponential-log} and~\ref{lem:twisted-obstruction}.
\foundation{F8}Lemma~\ref{lem:primitive} proves the all-finite-pole criterion
for $q=r'+\mu r$.
\foundation{F9}Lemma~\ref{lem:formal} formally diagonalizes systems with
distinct leading eigenvalues and proves invariance of powers modulo $\Z$.
\foundation{F10}For $a_2y''+a_1y'+a_0y=0$, differentiating the Wronskian
gives $W'=-(a_1/a_2)W$.
\foundation{F11}Sums, products, derivatives, algebraic rescalings and
primitives vanishing at zero preserve $E$-functions. The exponential,
rational-parameter Kummer and ${}_0F_1$ instances satisfy the exact bounds
and equations in Theorem~\ref{thm:values}'s proof.
\foundation{F12}We use DLMF 10.29.1(a), 10.29.3(a)--(b), 10.39.5, 10.40.1, 13.2.9, 13.2.39,
13.2.41, 13.3.17, 13.6.10 and 13.7.2--13.7.3 with their stated branches
and sectors. Our equations also imply the linkage maps.
\foundation{F13}Tensor products of geometrically integral algebras over a
field are geometrically integral \cite[Proposition~(4.6.5)(ii), p.~69]{EGAIV2}.
For a finitely generated integral $k$-algebra $A$, $\dim A=\trdeg_k A$
\cite[Theorem~5.6, p.~34]{Matsumura1989}. Consequently, for primes
$P\subset P'$ in a finitely generated $k$-algebra, equal quotient
dimensions imply $P=P'$.
We also use that free modules over domains are torsion-free and that
field extensions are free as modules.
\foundation{F14}Hermite normal form, number-field factorization and
Gr\"obner elimination are effective over the specified fields
\cite[\S\S2.4,3.6]{Cohen1993}, \cite{BeckerWeispfenning1993}.
Appendix~\ref{app:algorithm} refers to Supplementary Section~\ref{sup:input-model} for model (M-gen).
Its oracle (O1) gives integer-relation $\Z$-bases for finite tuples.
\foundation{F15}Algebraic-group morphisms have closed image; $G^0$ is normal of finite index.
Surjections preserve identity components
\cite[Proposition~1.2(b), p.~46; Corollary~1.4(a)--(b), p.~47]{Borel1991}.
In characteristic zero, $\dim G=\dim\Lie G$. This uses the smooth reduced-group
convention of \cite[\S9.1, p.~65]{Humphreys1975} together with the dimension statement in
\cite[Proposition~1.2(a), p.~46]{Borel1991}.
Also $\Lie\ker\phi=\ker d\phi$ and $\Lie(H\cap K)=\Lie H\cap\Lie K$
\cite[\S12.5, p.~85]{Humphreys1975};
$\Lie(\phi G)=d\phi(\Lie G)$ uses
\cite[\S7.2, proof, p.~105]{Borel1991}.
For connected $H$, $H\subset K$ iff $\Lie H\subset\Lie K$
\cite[\S13.1, Theorem, p.~87]{Humphreys1975}
\cite[\S7.1(2), p.~105]{Borel1991}.
For closed $H\subset G$, connected $G$ and $\dim H=\dim G$ imply $H=G$.
Apply proper-closed-subset dimension drop to $H^0$ in irreducible $G$
\cite[\S3.2, Proposition, p.~25]{Humphreys1975}
\cite[Proposition~1.2 and following paragraph, \mbox{pp.}~46--47]{Borel1991}.
Proper connected $\SL_2$ subgroups are solvable by
\cite[Theorem~4.29, pp.~133--134]{VanDerPutSinger2002Draft}.
In characteristic zero, nonidentity unipotent powers are dense in a connected
additive subgroup \cite[\S2.1, p.~56; \S7.3, Remark, p.~106]{Borel1991}.


For (F4), work over the algebraically closed characteristic-zero field
$\C$. Put $h=\operatorname{diag}(1,-1)$, $e=E_{12}$ and $f=E_{21}$.
For a Lie automorphism $\varphi$, $\operatorname{ad}\varphi(h)$ has
eigenvalues $0,2,-2$. A traceless $2\times2$ matrix with eigenvalues
$\lambda,-\lambda$ has adjoint eigenvalues $0,2\lambda,-2\lambda$. A nilpotent matrix has nilpotent adjoint. Thus $\varphi(h)$ has
eigenvalues $1,-1$. Inner conjugation makes $\varphi(h)=h$.
The weight spaces give $\varphi(e)=ae$, $\varphi(f)=bf$. Brackets give $ab=1$. Choosing $t\in\C^\times$ with $t^2=a$ identifies the
remaining automorphism as conjugation by $\operatorname{diag}(t,t^{-1})$.

For (F15), let $G\subseteq\operatorname{GL}_2(\C)$ be a connected reductive
algebraic subgroup in its defining representation, and put $H=G\cap\SL_2$.
If $H^0$ is proper, it lies in a Borel subgroup or torus and is solvable.
The connected commutator $[G,G]$ lies in $H^0$. Since $G/H$ is abelian,
Corollary~A.38 makes $G$ solvable. Lie--Kolchin and complete reducibility
give an invariant line and complement. Thus $G$ is conjugate to a connected
diagonal subgroup, hence a torus. A non-torus $G$ therefore contains $\SL_2$.
The solvability and invariant-line inputs
are \cite[Corollary~A.38, p.~364; Theorem~A.46, p.~367]{VanDerPutSinger2003}.
Complete reducibility is stated in the unnumbered paragraph after
Definition~2.37, p.~59, of the July 2002 draft
\cite{VanDerPutSinger2002Draft}.

\subsection{Formal comparison}\label{app:formal}
\begin{lemma}\label{lem:formal}
For rank-two $A=A_0+O(z^{-1})$ with distinct leading eigenvalues,
rational isomorphisms preserve those eigenvalues and formal powers modulo
$\Z$, and conjugate finite local monodromy.
\end{lemma}
\begin{proof}
In an eigenbasis solve $AP-P'=PB$, with $P=1+\sum P_nz^{-n}$ and
diagonal $B$, recursively by
\begin{equation}\label{eq:core-formal-recursion}
\begin{gathered}
R_n=A_n+\sum_{j=1}^{n-1}(A_jP_{n-j}-P_{n-j}B_j)+(n-1)P_{n-1},\\
B_n=\operatorname{diag}R_n,\quad(P_n)_{ij}=-(R_n)_{ij}/(\omega_i-\omega_j).
\end{gathered}
\end{equation}
Take $P_n$ with zero diagonal. A nonzero intertwiner entry has logarithmic
derivative $\omega_i-\omega_j+(\vartheta_i-\vartheta_j)/z+O(z^{-2})$,
versus $n/z+O(z^{-2})$ for a Laurent series. Invertibility pairs entries. Rational gauges are single-valued at finite points.
\end{proof}

\subsection{Arithmetic coefficient bounds}\label{app:arithmetic}
For the $E$-function bounds, write each coordinate as $\sum a_nz^n/n!$,
with coefficients in one number field and conjugates of modulus $e^{O(n)}$.
For an algebraic tuple $a$, $\operatorname{den}(a)$ is the least positive
integer $m$ making every entry of $ma$ an algebraic integer. For
$G=\sum g_jz^j/j!$, put $G_{\le m}=(g_0,\ldots,g_m)$.
Put $L_n=\operatorname{lcm}(1,\ldots,n)=e^{O(n)}$. Common denominators
through the indicated degree divide
\[
\begin{array}{c|c}
\Ein_k(\eta z),\ n&\operatorname{den}(\eta)^nL_n^k\\
I_c,S_c,\ 2n&4^n\operatorname{den}(c)^{2n}L_n\\
z^{-N}G,\ n\ (G=O(z^N))&\operatorname{den}(G_{\le n+N})L_{n+N}^N.
\end{array}
\]
For the order-zero Bessel coordinates, the coefficients are $\binom{2m}m(c/2)^{2m}$ and its product
with $\mathsf H_m$. The third table row also gives the $(\mathrm{B2}^E)$ bounds by zero-origin
integration. For Kummer $a=p/q$ in lowest terms with $q>0$, $(a)_m/m!$ has denominator dividing
$q^{2m}$. Primes not dividing $q$ cancel $v_\ell(m!)$ in the progression
$p+jq$, and $v_\ell(m!)<m$ for $\ell\mid q$. The denominator of
$h_m=2\mathsf H_m-\sum_{j<m}q/(p+jq)$ divides $L_{qm+|p|}$.
For nonresonant Bessel, with $\nu+1=p/q$, $\gcd(p,q)=1$ and $q>0$,
\begin{equation}\label{eq:bessel-normalized-coefficient}
a_{2m}=c^{2m}(1/2)_m/(\nu+1)_m,\qquad a_{2m+1}=0.
\end{equation}
Let $d_m$ be the denominator of $(1/2)_m/(\nu+1)_m$ in lowest terms.
For odd $\ell\nmid q$, denominator and odd-numerator progression counts
differ by at most one at each $\ell^e$. Odd $\ell\mid q$ give no denominator.
At $2$, $2^{-m}$ and $\sum\lfloor m/2^e\rfloor\le m$ give
\begin{equation}\label{eq:bessel-common-denominator}
d_m\mid4^mL_{|p|+qm},\quad
\operatorname{den}(a_0,\ldots,a_{2n})\mid
\operatorname{den}(c)^{2n}4^nL_{|p|+qn}.
\end{equation}
Replacing $\nu$ by $-\nu$ gives the other solution. The Pochhammer ratio grows polynomially. This gives the conjugate bounds. Differentiation shifts $m$ to $m+1$ in the Kummer coefficient formulas
and preserves the $e^{O(n)}$ conjugate and denominator bounds.


\section{The input model and computation}\label{app:algorithm}
Supplementary Section~\ref{sup:input-model} gives the full input
model, computability proposition and proof.
It requires exact presented fields and compatible normalization constants.
Approximations alone are excluded.
Appendix~\ref{app:arithmetic} supplies the value coefficient bounds.

\section{Parameters in Gorelov's classification}\label{app:gorelov}
Gorelov uses data $(l,q,p;\boldsymbol\nu,\boldsymbol\lambda,\alpha)$.
His 2021 Lemmas~12 and~14 treat $p_1=p_2=1$ and arbitrary powers,
respectively (statements p.~7; proofs p.~7 and \mbox{pp.}~8--9)
\cite[Lemmas~12 and~14]{Gorelov2021}.
The proof of Lemma~12 in \cite[p.~7]{Gorelov2021} invokes
Lemma~14 of that article's reference~[5] (Gorelov, \emph{Math. Notes}
94 (2013), 82--95), then Lemmas~1--2 of the citing article.
His 2020 Theorem~2 gives mixed sufficiency
\cite[p.~138]{Gorelov2020}, extending the 2019 identity
\cite[(3), p.~260]{Gorelov2019}. The 2019 matrix form
\cite[Theorem~1, p.~260]{Gorelov2019} assumes $2\lambda\notin\Z$
and excludes our case $\nu_0=0$.
Theorem~\ref{thm:linkage} uses
\begin{equation*}
\begin{gathered}
\begin{array}{c|ccc|c|c|c}
 &l&q&p&\boldsymbol\nu&\boldsymbol\lambda&\alpha\\ \hline
 \operatorname{Bess}(c,\nu_0)&0&2&2&\varnothing&(1,1+\nu_0)&c^2/4\\
 \operatorname{Kum}(s,a)&1&2&1&(a)&(1,1)&s
\end{array}
\end{gathered}
\end{equation*}
BBH Theorem~2.2(iv), p.~288, applies to Lemma~\ref{lem:groups}'s scalar
extensions of $\SL_2$ \cite{BBH1988}. It supplies irreducible Siegel-normal
homogeneous systems for Proposition~2.4, \mbox{pp.}~291--292, including logarithmic
second solutions but excluding separately adjoined $L$ and E-parts;
see Supplementary Section~\ref{sup:classical-comparison}. Admissibility conditions: \mbox{pp.}~280--281, 298--299;
rational-value conclusion: Corollary~4.6, \mbox{pp.}~300--301.
Gorelov records the Bessel--Kummer omission and corrected homogeneous condition
\cite[Example~1 and (9), p.~3; Lemma~14, \mbox{pp.}~7--9]{Gorelov2021}.

Scales are nonzero. Denominator parameters avoid $\Z_{\le0}$. Lemma~\ref{lem:groups} gives irreducibility. The formal data are
\begin{equation*}
\begin{gathered}
\begin{array}{c|c|c|c}
 \text{block}&\text{Wronskian, up to a constant}&
   (\text{rate},\text{power})\text{ at }\infty&\text{exponents at }0\\ \hline
 \operatorname{Bess}(c,0)&z^{-1}&(\pm c,-1/2)&0,0\text{ with log}\\
 \operatorname{Bess}(c,\nu)&z^{-2\nu-1}&(\pm c,-\nu-1/2)&0,-2\nu\\
 \operatorname{Kum}(s,a)&e^{sz}/z&(0,-a),(s,a-1)&0,0\text{ with log}
\end{array}
\end{gathered}
\end{equation*}
The powers follow from \cite[10.40.1, 13.7.2--13.7.3]{DLMF} and are
preserved modulo $\Z$ by Lemma~\ref{lem:formal}. At zero, distinct
nonresonant eigenvalues versus nontrivial resonant unipotence give
Theorem~\ref{thm:linkage}'s exclusions.

For the Bessel--Bessel necessity, use the convention
$M_j\simeq M_i\otimes e^{\lambda z}z^\upsilon$.
The rates at infinity give
$\{c',-c'\}=\{c+\lambda,-c+\lambda\}$. Thus $\lambda=0$ and
$c'=\pm c$. The powers then give $\upsilon\equiv\nu_0-\nu'_0\pmod\Z$.
At zero, matching the exponent multisets requires
$\{0,-2\nu'_0\}\equiv\{\upsilon,\upsilon-2\nu_0\}\pmod\Z$.
If $\upsilon\equiv0$, the infinity congruence gives $\nu'_0\equiv\nu_0$. If $\upsilon\equiv2\nu_0$, it gives $\nu'_0\equiv-\nu_0$.
In $[0,1/2)$ the first implies equality and the second forces both zero. Hence $\nu'_0=\nu_0$.

\section{Proof of the explicit linear independence estimate}\label{sec:measure}

\begin{explicitestimate}[Evaluated constants]\label{thm:measure}
Put
\begin{equation}\label{m4:constants}
 A_*=311040\cdot2^{24001},\qquad N=A_*^8,\qquad H_0=N^N.
\end{equation}
For $(a,b,c,d)\in\Z^4\setminus\{0\}$ and
$H=\max(|a|,|b|,|c|,|d|)\ge H_0$,
\begin{equation}\label{m4:measure}
 |a+be+ce^{-1}+d\Ein(1)|\ge\frac1{12}H^{-60}.
\end{equation}
\end{explicitestimate}

We evaluate the construction of
\cite[Proposition~3.1 and Lemma~3.2, pp.~12--15]{AdamczewskiFaverjon2026}
at $\nu=\mu=4$. Below we prove linear independence over $\C(z)$
and select rows by direct Wronskian cyclicity, replacing a generic zero estimate.

\begin{lemma}\label{m4:matrix}
For every integer $n\ge2^{16}$, the finite search and exact row selection below construct
$Q_n\in M_4(\Z)$, with
\begin{equation}\label{m4:matrix-bounds}
 \det Q_n\ne0,\qquad
 \max_{i,k}|(Q_n)_{i,k}|\le A_*^nn^{9n/8},\qquad
 \|Q_nf(1)\|_\infty\le A_*^nn^{-23n/8},
\end{equation}
where $f=(1,e^z,e^{-z},\Ein_1(z))^t$.
\end{lemma}

\begin{proof}
Write $f=(1,u,w,v)^t$, $D=d/dz$ and $\vartheta=zD$.
The system and polynomial-row update are
\begin{equation}\label{m4:system}
 u'=u,\qquad w'=-w,\qquad v'=(1-w)/z,
\end{equation}
\begin{equation}\label{m4:row-update}
 \vartheta(Pf)=\mathscr U(P)f,\quad
 \mathscr U(P)=
 (zP_0'+P_v,\ zP_+'+zP_+,\ zP_-'-zP_--P_v,\ zP_v').
\end{equation}
Clearing denominators and comparing growth at $\pm\infty$ proves
independence of $1,u,w$ over $\C(z)$. If
$v=r_0+r_+u+r_-w$ with rational $r_0,r_+,r_-$, differentiating and
comparing constant coefficients gives $r_0'=1/z$, contradicting the
zero residue of rational derivatives at $0$. Thus $1,u,w,v$ are independent.

\emph{Integer polynomials.}
Fix $n\ge2^{16}$, and put
\begin{equation}\label{m4:parameters}
\begin{gathered}
 N_{\rm var}=4(n+1),\qquad m=4n+3-\lfloor n/240\rfloor,\\
 C_4=2^{24001},\qquad B_n=(3N_{\rm var}\,8^{5n})^{1200}+1.
\end{gathered}
\end{equation}
Write $f_i(z)=\sum_{t\ge0}a_{i,t}z^t/t!$.
Here $a_{0,t}=\mathbf1_{t=0}$, $a_{+,t}=1$,
$a_{-,t}=(-1)^t$, $a_{v,0}=0$, and
$a_{v,t}=(-1)^{t+1}/t$ for $t\ge1$.
Thus $|a_{i,t}|\le1$.
With $L_m=\operatorname{lcm}(1,\ldots,m)$ impose
\begin{equation}\label{m4:integer-system}
 \sum_{i\in\{0,+,-,v\}}\sum_{j=0}^{\min(n,t)}
 L_m\binom tj a_{i,t-j}p_{i,j}=0\qquad(0\le t<m).
\end{equation}
This gives $m$ integer equations in $N_{\rm var}$ unknowns. To bound their
coefficients, we use $L_s\le4^s$, which follows by induction from
\[
 L_{2r}/L_r\mid\binom{2r}r,\qquad
 L_{2r+1}/L_{r+1}\mid\binom{2r+1}r.
\]
The factorial valuation formula gives these divisibilities. Each new maximal prime power lies in the upper half. Induct from $L_1=1$ using
$\binom{2r}r\le2^{2r}$ and $\binom{2r+1}r\le2^{2r}$.
Each entry of \eqref{m4:integer-system} is therefore at most $8^m\le8^{5n}$.

The integer image of $\{0,\ldots,B_n\}^{N_{\rm var}}$ has at most
$(3N_{\rm var}8^{5n}B_n)^m$ elements. Also
\[
 N_{\rm var}-m>n/240,\qquad \frac m{N_{\rm var}-m}<1200,\qquad
 B_n^{N_{\rm var}-m}>(3N_{\rm var}8^{5n})^m.
\]
There are more box vectors than images.
Two distinct vectors have the same image.
Their difference gives a nonzero solution with
\begin{equation}\label{m4:box-bound}
 \max|p_{i,j}|\le B_n\le2^{24000n}+1\le C_4^n,
 \qquad 3N_{\rm var}\le15n\le2^{5n}.
\end{equation}
The image bound makes the lexicographic search of
$[-B_n,B_n]^{N_{\rm var}}\cap\Z^{N_{\rm var}}$, excluding zero, terminate at a
solution. Set
\begin{equation}\label{m4:polynomials}
 P_i=n!\sum_{j=0}^np_{i,j}z^j/j!\in\Z[z],\qquad
 R=Pf\not\equiv0.
\end{equation}
Then $\deg P_i\le n$, $h(P_i)\le C_4^nn!$, and $\operatorname{ord}_0R\ge m$,
where $h$ denotes maximum coefficient modulus.

\emph{Cyclicity.}
If $P_v=0$, the nonzero $R$ satisfies the monic constant-coefficient equation
\[
 D^{n+1}(D-1)^{n+1}(D+1)^{n+1}R=0.
\]
Its order-$3n+3$ Taylor recurrence gives
$\operatorname{ord}_0R\le3n+2<m$, a contradiction. Thus $P_v\ne0$. Normalize $r=R/P_v=v+r_0+r_+u+r_-w$. Then
\begin{equation}\label{m4:derivative}
 \begin{aligned}
 r'&=q_0+B_-w+B_+u,\qquad q_0=z^{-1}+r_0',\\
 B_-&=-z^{-1}+r_-'-r_-,\qquad B_+=r_+'+r_+.
\end{aligned}
\end{equation}
The residue at $0$ gives $q_0\ne0$. No rational function solves
$r_-'-r_-=1/z$: differentiation increases any finite pole order, and
polynomials cannot produce $1/z$. Thus $B_-\ne0$. The same pole argument
and degree comparison give $B_+\ne0$ whenever $r_+\ne0$.

For meromorphic functions, Wronskian vanishing is equivalent to constant
dependence, by induction using
\[
 W(h_1,\ldots,h_k)
 =h_1^kW\bigl((h_2/h_1)',\ldots,(h_k/h_1)'\bigr).
\]
Work locally where $h_1\ne0$, then continue meromorphically.
If $P_+=0$, independence of $q_0,B_-w$ gives
$\det\operatorname{coeff}_{(1,w)}(r',r'')=W(q_0,B_-w)/w\ne0$.
Adding $r$ gives rank three on $(1,w,v)$.
Multiplication by $P_v$ and replacement of $D^j$ by $\vartheta^j$
are row operations over $\C(z)$. Both are invertible and triangular.
Hence the coefficient determinant
$\Delta_3$ of $R,\vartheta R,\vartheta^2R$ on $(1,w,v)$ is a nonzero polynomial
of degree at most $3n+3$.
Its adjugate, applied to the constant coordinate, gives
$\Delta_3=\sum_{j=0}^2c_j\vartheta^jR$, with polynomial $c_j$.
Hence $\operatorname{ord}_0\Delta_3\ge m>3n+3$, a contradiction.
Therefore $P_+\ne0$.
Now $q_0,B_-w,B_+u$ are independent over $\C$.
Their Wronskian gives
\[
 \det\operatorname{coeff}_{(1,w,u)}(r',r'',r''')
 =\frac{W(q_0,B_-w,B_+u)}{wu}\ne0.
\]
Adding $r$ and repeating the triangular operations gives independent
coefficient rows for $R,\vartheta R,\vartheta^2R,\vartheta^3R$.

\emph{Selection at $1$.}
Put
\[
 \ell_0=P,\qquad \ell_{j+1}=\mathscr U(\ell_j),\qquad R_j=\ell_jf,\qquad
 \Delta=\det(\ell_0,\ldots,\ell_3)\ne0.
\]
Each $R_j=\vartheta^jR$ has order at least $m$ at zero.
The adjugate expresses $\Delta$ in their polynomial span, giving
\begin{equation}\label{m4:determinant-orders}
 \deg\Delta\le4n+6,\quad \operatorname{ord}_0\Delta\ge m,\quad
 \chi:=\operatorname{ord}_1\Delta\le3+\lfloor n/240\rfloor.
\end{equation}
On $Y=(Y_0,Y_+,Y_-,Y_v)^t$, extend $\vartheta$ to
\[
 \mathcal D=z\partial_z+zY_+\partial_{Y_+}-zY_-\partial_{Y_-}
 +(Y_0-Y_-)\partial_{Y_v}.
\]
Since $\mathcal D(\ell_jY)=\ell_{j+1}Y$, apply $\mathcal D^\chi$
to the adjugate identity $\Delta Y_i=\sum_{j=0}^3c_{i,j}\ell_jY$.
At $1$, Leibniz's rule gives
\begin{equation}\label{m4:formal-selection}
 (\vartheta^\chi\Delta)(1)Y_i
 =\sum_{j=0}^3\sum_{s=0}^\chi\binom\chi s
 (\vartheta^{\chi-s}c_{i,j})(1)\ell_{j+s}(1)Y.
\end{equation}
Lower derivatives vanish, but $(\vartheta^\chi\Delta)(1)=D^\chi\Delta(1)\ne0$.
Hence $\ell_j(1)$, $0\le j\le J:=6+\lfloor n/240\rfloor$, span $\Q^4$.
Gaussian elimination chooses their first invertible four-row subset as $Q_n$.

\emph{Heights and tails.}
For $j\le J$, row degrees are at most $n+j\le2n$.
Update \eqref{m4:row-update} and evaluation at $1$ each cost at most $3n$ in height.
Using $j+1\le n/240+7\le n/8$ gives
\begin{equation}\label{m4:height}
 \max_k|\ell_{j,k}(1)|\le(3n)^{j+1}C_4^nn!
 \le(3C_4)^nn^{9n/8}.
\end{equation}
Write $R=\sum_{t\ge m}r_tz^t$.
For $t\ge m\ge3n$ and $j\le J$, convolution and \eqref{m4:box-bound}
give the coefficient bound and consecutive tail ratio
\[
\begin{gathered}
 |r_t|\le\frac{4n!C_4^n}{t!}\sum_{j=0}^{\min(n,t)}\binom tj
 \le4n!C_4^n\frac{2^t}{t!},\\
 \frac2{t+1}(1+1/t)^j\le\frac{2e^{j/t}}{t+1}
 \le\frac6{t+1}\le\frac12.
\end{gathered}
\]
The geometric tail and $m!\ge(m/3)^m$ now give
\begin{align}
 |R_j(1)|&\le4n!C_4^n\sum_{t\ge m}\frac{t^j2^t}{t!}
 \le8n!C_4^n\frac{m^j2^m}{m!}\notag\\
 &\le8C_4^n6^{5n}n^{n+j-m}\notag\\
 &\le(40\cdot6^5C_4)^nn^{-23n/8}.\label{m4:tail-final}
\end{align}
Here $\log(m!)\ge\int_1^m\log x\,dx$ justifies the factorial bound.
Since $n\le m\le5n$ and $j\le n\le m$, use $m^{j-m}\le n^{j-m}$ in the second inequality.
For $n\ge26$, $n+j-m\le-3n+n/120+3\le-23n/8$.
Since $40\cdot6^5C_4=A_*\ge3C_4$, the height and tail bounds give \eqref{m4:matrix-bounds}.
\end{proof}

\begin{proof}[Proof of Theorem~\ref{thm:measure}]
The least $n\ge N$ with $n^n\ge H^8$ satisfies $n>N$ and
$(n-1)^{n-1}<H^8$, since $H\ge N^N$.
For $n\ge4$, $(n/(n-1))^{n-1}<3$ and $(n-1)^{n-1}\ge3n$ give
\begin{equation}\label{m4:height-transfer}
 H\le n^{n/8},\qquad n^n\le(n-1)^{2(n-1)}<H^{16}.
\end{equation}
As $n\ge N=A_*^8>2^{16}$, $A_*\le n^{1/8}$.
Lemma~\ref{m4:matrix} gives $B^*_n=A_*^nn^{9n/8}\le n^{5n/4}$
and $E_n=A_*^nn^{-23n/8}\le n^{-11n/4}$.
Three rows of invertible $Q_n$ complete $p=(a,b,c,d)\ne0$
to an invertible integer matrix $M$ with first row $p$.
Thus $|\det M|\ge1$.
Put $\Lambda=pf(1)$. Cofactors have six terms, giving $6(B^*_n)^3$
for $\Lambda$ and $6H(B^*_n)^2$ per other row.
The first coordinate of $\operatorname{adj}(M)Mf(1)=(\det M)f(1)$ uses $f_0(1)=1$.
Since $n\ge64$, $18n^{-n/8}<1/2$.
Hence \eqref{m4:height-transfer} yields
\begin{align}
 1&\le6(B^*_n)^3|\Lambda|+18H(B^*_n)^2E_n
 \le6(B^*_n)^3|\Lambda|+18n^{-n/8},\label{m4:cofactor}\\
 |\Lambda|&\ge\frac1{12(B^*_n)^3}\ge\frac1{12}n^{-15n/4}
 >\frac1{12}H^{-60}.\notag\qedhere
\end{align}
\end{proof}

\section{Detailed proofs of the Galois-group and specialization steps}
\label{app:detailed-proofs}

\subsection{The full joint group}\label{app:joint-group}
We expand the group argument in Proposition~\ref{prop:product}.
All groups are reduced complex linear algebraic groups.

\begin{lemma}\label{lem:app-group-interface}
Let $\phi:G\to H$ be a group morphism.
Its image is closed, $G^0$ is normal of finite index, and
\begin{equation}\label{eq:app-group-interface}
\begin{gathered}
 \phi(G^0)=\phi(G)^0,\qquad \dim G=\dim\Lie G,\\
 \Lie(\ker\phi)=\ker d\phi,\qquad
 \Lie(\phi(G))=d\phi(\Lie G),\\
 \Lie(A\cap B)=\Lie A\cap\Lie B
 \quad(A,B\subset G\text{ closed}).
\end{gathered}
\end{equation}
If $A$ is connected, then $\Lie A\subset\Lie B$ implies $A\subset B$.
A closed subgroup of a connected group with equal dimension is the whole group.
\end{lemma}
\begin{proof}
Closed images and identity components follow from
\cite[Proposition~1.2(b), p.~46; Corollary~1.4(a)--(b), p.~47]{Borel1991}.
Smoothness identifies group dimension with tangent-space dimension
\cite[Proposition~1.2(a), p.~46]{Borel1991}
\cite[\S9.1, p.~65]{Humphreys1975}.
The kernel and intersection identities are
\cite[\S12.5, Theorem, p.~85]{Humphreys1975}.
The image identity uses characteristic-zero separability
\cite[\S7.2, proof, p.~105]{Borel1991}.
If $\Lie A\subset\Lie B$, then $A\cap B$ has dimension $\dim A$.
Connected algebraic groups are irreducible
\cite[Proposition~1.2 and following paragraph, \mbox{pp.}~46--47]{Borel1991}.
Proper closed irreducible subsets have smaller dimension
\cite[\S3.2, Proposition, p.~25]{Humphreys1975}.
Apply this to $(A\cap B)^0\subset A$.
It gives $A\subset B$; see also
\cite[\S7.1(2), p.~105]{Borel1991}
\cite[\S13.1, Theorem, p.~87]{Humphreys1975}.
The last assertion follows by applying the same argument to $B^0\subset G$.
\end{proof}

\begin{lemma}\label{lem:app-pv-compositum}
Let $E/K$ and $F/K$ be PV extensions with algebraically closed constants $\C$.
Assume their compositum $M=EF$ has the same constants.
Then $M/K$ and $M/F$ are PV, and restriction gives
\begin{equation}\label{eq:app-pv-restriction}
 \Gal(M/F)\simeq\Gal(E/(E\cap F)).
\end{equation}
The right side is normal in $\Gal(E/K)$.
Its quotient is a quotient of $\Gal(F/K)$.
The intersection $E\cap F$ is PV over $K$.
The maps $\Gal(M/K)\to\Gal(E/K),\Gal(F/K)$ are onto.
\end{lemma}
\begin{proof}
The block sum of fundamental matrices generates $M/K$.
The matrix for $E$ also generates $M/F$.
Both subfields $E,F$ are invariant under $\Gal(M/K)$.
PV Galois correspondence gives the two surjections
\cite[Proposition~1.34(1)--(2), p.~26]{VanDerPutSinger2002Draft}.
The restriction $\Gal(M/F)\to\Gal(E/K)$ is injective: an automorphism
in its kernel fixes $E$ and $F$, hence fixes $M=EF$, so it is the identity.
Put $H=\operatorname{im}(\Gal(M/F)\to\Gal(E/K))$.
It is closed by Lemma~\ref{lem:app-group-interface}.
Surjectivity and conjugation show that $H$ is normal.
Its fixed field is $E\cap F$, proving
\eqref{eq:app-pv-restriction} by the same correspondence.
Since $H$ is normal, the PV correspondence also makes $E\cap F$ PV over $K$
\cite[Proposition~1.34(2), p.~26]{VanDerPutSinger2002Draft}.
The map to $\Gal(E/K)/H$ kills $\Gal(M/F)$.
It therefore factors through $\Gal(F/K)$.
\end{proof}

\begin{proof}[Detailed proof of Proposition~\ref{prop:product}]
\emph{1. Base change.}
Put $V=\mathsf G_{u,\mathcal U}$ and $H_0=\Gal(QT'/T')$.
The matrices in \eqref{eq:core-unipotent-matrix} are lower unitriangular.
Thus $V$ is a closed unitriangular subgroup, hence nilpotent.
Every nonidentity unipotent has connected cyclic closure
\cite[\S2.1, p.~56; \S7.3, Remark, p.~106]{Borel1991}.
Each element therefore belongs to $V^0$.
Hence $V$ is connected and solvable.

The generators of $T'$ satisfy
\[
 t_h'=\lambda_ht_h,\qquad
 (z^{1/\ell})'=\frac{z^{1/\ell}}{\ell z}.
\]
Thus $T'/k$ is PV with abelian group.
For each block field $\mathcal P_i$, Lemma~\ref{lem:app-pv-compositum}
identifies $\Gal(\mathcal P_iT'/T')$ with a normal subgroup of
$\Gal(\mathcal P_i/k)$ having abelian quotient.
The block group contains $\SL_2$ by Lemma~\ref{lem:groups}.
Since $\SL_2=[\SL_2,\SL_2]$, the normal subgroup contains $\SL_2$.
The group over $T'$ fixes the block Wronskian, which belongs to $T'$:
\[
 W_i\in\C^\times\{z^{-1},\ z^{-2\nu-1},\ e^{sz}/z\}.
\]
Thus the base-changed block group equals $\SL_2$.
All composita have constants $\C$, by the standing constant-field hypothesis.

\emph{2. Projective independence.}
From \eqref{eq:core-block-transport},
\[
 w_i\mathbf1=Y_{\mathfrak c}^{-1}B_i^{-1}Y_iC_i^{-1}
 \quad\Longrightarrow\quad w_i\in Q\cap T'.
\]
Thus \eqref{eq:core-character} holds on the full group $G_Q$.
It gives $\psi(G_Q)\subset\bar P_{\mathrm{SL}}=\psi(P_{\mathrm{SL}})$,
where $P_{\mathrm{SL}}$ is the subgroup denoted $P$ in the body, with its conjugated block embedding.
Each representative projection is onto $\PGL_2$.
Each block field is a PV subextension of $Q/k$.
Lemma~\ref{lem:app-group-interface} makes
$\Lie(\psi(G_Q)^0)$ subdirect in the representative factors.

We use an elementary consequence of Lemma~\ref{lem:goursat}.
A proper subdirect Lie algebra in $\mathfrak{sl}_2^r$ has a graph projection
onto some pair of factors.
For $r=2$, use Lemma~\ref{lem:goursat}.
Proceed by induction on $r$.
If the first $r-1$ projection is proper, apply induction.
Otherwise it is $\mathfrak s=\mathfrak{sl}_2^{r-1}$.
For $\mathfrak h\subset\mathfrak s\oplus\mathfrak{sl}_2$, put
$\mathfrak j=\{b:(0,b)\in\mathfrak h\}$.
It is an ideal in $\mathfrak{sl}_2$.
If $\mathfrak j=\mathfrak{sl}_2$, subdirectness gives the full sum.
Otherwise $\mathfrak j=0$, and $\mathfrak h$ is the graph of a surjection
$\mathfrak s\to\mathfrak{sl}_2$.
The factor images are commuting ideals.
Exactly one is nonzero, giving the required pair.
The standard brackets $[h,e]=2e$, $[h,f]=-2f$, and $[e,f]=h$
show that $\mathfrak{sl}_2$ is simple.
The calculation in (F4) shows that its automorphisms are inner.

A graph pair corresponds to a rank-one twist by Lemma~\ref{lem:twist}.
For completeness, let $G_Q=\Gal(Q/k)$ be the full joint group used in descent.
The projective image of $G_Q^0$ in the $i,j$ factors is the graph group by
Lemma~\ref{lem:app-group-interface}.
Thus $\pi_j=\chi C_0\pi_iC_0^{-1}$ on $G_Q^0$.
On $\operatorname{Hom}(V_i,V_j)$, the line $\C C_0$ is the unique invariant line:
the other summand is the irreducible adjoint representation.
Its invariant subspaces are ideals in $\mathfrak{sl}_2$.
Normality of $G_Q^0$ makes this line $G_Q$-stable.
PV descent gives a rational basis $B_0$ of that line.
In the Hom system it has the form
\[
 B_0=vY_jC_0Y_i^{-1},\qquad
 B_0'+B_0M_i-M_jB_0=rB_0,\qquad r=v'/v\in k.
\]
Here $C_0\in\GL_2(\C)$ is the graph matrix, and $M_i,M_j$ are the block
system matrices.
Thus $B_0$ is invertible.
A scalar solution of $u'/u=-r$ turns it into an intertwiner.
Distinct representatives belong to distinct linkage classes.
They therefore have no graph pair.
The induction gives the full direct-sum Lie algebra.
The dimension criterion in Lemma~\ref{lem:app-group-interface} now yields
\begin{equation}\label{eq:app-projective-full}
 \psi(G_Q^0)=\psi(G_Q)=\bar P_{\mathrm{SL}}.
\end{equation}

\emph{3. The full group after base change.}
Wronskians and \eqref{eq:core-character} give $H_0\subset P_{\mathrm{SL}}$.
Restriction $\Gal(QT'/k)\to G_Q$ is onto.
Lift any $g,h\in G_Q$ to this group.
Their restrictions to $T'$ commute.
Their commutator therefore fixes $T'$, giving $[g,h]\in H_0$.
For $x,y\in P_{\mathrm{SL}}$, equation~\eqref{eq:app-projective-full} supplies $g,h\in G_Q$
with the same respective projective images.
Thus $g=xz_1$, $h=yz_2$, where $z_1,z_2\in\ker\psi$ are central.
Hence $[x,y]=[g,h]\in H_0$.

Put $u_+(b)=\left(\begin{smallmatrix}1&b\\0&1\end{smallmatrix}\right)$,
$u_-(b)=\left(\begin{smallmatrix}1&0\\b&1\end{smallmatrix}\right)$ and
$d(t)=\operatorname{diag}(t,t^{-1})$.
Use the commutator convention $[x,y]=xyx^{-1}y^{-1}$.
For $t^2\ne1$,
\[
 u_+(b)=[d(t),u_+(b/(t^2-1))],\qquad
 u_-(b)=[d(t),u_-(b/(t^{-2}-1))].
\]
Row elimination shows that these elementary matrices generate $\SL_2(\C)$.
Thus $P_{\mathrm{SL}}$ is perfect, and $P_{\mathrm{SL}}\subset H_0$.
Consequently $H_0=P_{\mathrm{SL}}$, including all finite components.

\emph{4. The unipotent factor.}
Both $\mathcal U/T'$ and $QT'/T'$ are PV.
Lemma~\ref{lem:app-pv-compositum} gives surjective restriction maps from
$\Gal(\Omega/T')$ to both groups.
Their product is injective since $\Omega=\mathcal U(QT')$.
Write its image as $H\subset V\times P_{\mathrm{SL}}$ and put
\[
 H_V=\{v:(v,1)\in H\},\qquad H_P=\{p:(1,p)\in H\}.
\]
Conjugation by lifts proves normality in $V,P_{\mathrm{SL}}$, respectively.
Multiplication of lifts gives the isomorphism
\[
 V/H_V\longrightarrow P_{\mathrm{SL}}/H_P,
 \qquad vH_V\longmapsto pH_P
 \quad((v,p)\in H).
\]
Two choices of $p$ differ by $H_P$. The map
$V\to P_{\mathrm{SL}}/H_P$ has kernel $H_V$, so it induces the displayed
isomorphism.
The common quotient is solvable and perfect, hence trivial.
Thus $H=V\times P_{\mathrm{SL}}$, proving \eqref{eq:core-product}.
Also $\mathcal U\cap QT'=T'$: its PV group is the same common quotient,
by Lemma~\ref{lem:app-pv-compositum}.
\end{proof}


\subsection{The tower induction}\label{app:tower-induction}
Keep the hypotheses of Lemma~\ref{lem:tower}.
The nodes $\eta_1,\ldots,\eta_m$ are distinct and nonzero.
The heights $h_j$ are positive integers.
Every rate in $\mathfrak S$ belongs to $\Gamma$.
The finite support contains $(0,0)$ and every $(-\eta_j,0)$.
Use $\Phi_{0,0}=L$ and $\Phi_{-\eta_j,0}=G_{\eta_j,1}$,
as in \eqref{eq:core-G}.
Put
\begin{equation}\label{eq:appendix-tower-assertions}
\begin{gathered}
 H(h)=\sum_j(h_j-1),\qquad
 \mathcal T_h(\mathbf p)=\sum_jp_j(L)G_{\eta_j,h_j}/z,\\
 \trdeg_T\mathcal N_h=\#\mathfrak S+H(h),\qquad
 \mathcal N_h^D=\C,\\
 \mathcal T_h(\mathbf p)\notin D\mathcal N_h
 \quad(0\ne\mathbf p\in\C[L]^m).
\end{gathered}
\end{equation}
The degree assertion means independence of all displayed generators.
We prove these assertions by induction on $H(h)$.

\begin{lemma}\label{lem:appendix-primitive-extension}
Let $K$ be a differential field of characteristic zero.
Suppose $K^D=\C$, $a\in K\setminus DK$ and $Du=a$.
Then $u$ is transcendental over $K$, and $K(u)^D=\C$.
\end{lemma}
\begin{proof}
If $u$ were algebraic, put $n=[K(u):K]$.
Derivation commutes with trace, as in Lemma~\ref{lem:primitive}.
Thus
\[
 D\bigl(n^{-1}\operatorname{Tr}_{K(u)/K}u\bigr)=a,
\]
contrary to $a\notin DK$.
For $Df=0$, write $f=P(u)/Q(u)$ in lowest terms.
Take $Q$ monic.
The identity $QDP=PDQ$ gives $Q\mid DQ$.
Monicity gives $\deg_u DQ<\deg_u Q$, so $DQ=0$.
If $\deg_u Q=r>0$, its next coefficient gives
$Dq_{r-1}+ra=0$.
Hence $Q=1$.
If $\deg_u P=s>0$, then $p_s\in\C^\times$ and
$Dp_{s-1}+sp_sa=0$.
Both equations contradict $a\notin DK$.
Thus $f\in\C$.
\end{proof}

For independent variables $Y_i$ over $K$, put $\partial_i=\partial/\partial Y_i$.
Their common kernel in $K(Y_i)$ is $K$.
A reduced numerator and denominator must divide their partial derivatives.
Degree comparison makes those derivatives zero.
If $DY_i\in K$, then $[D,\partial_i]=0$.
This follows on $K$ and on each $Y_i$.

\emph{1. Height zero.}
The vectors $\operatorname{res}(\Phi_\sigma)$ are nonzero multiples
of distinct unit vectors.
Lemma~\ref{lem:primitive} gives independence and constant field $\C$.
Put $\mathfrak U=\{(-\eta_j,0):1\le j\le m\}$ and
$\mathfrak S_0=\mathfrak S\setminus\{(0,0)\}$.
Suppose $Dy=\mathcal T_{\mathbf1}(\mathbf p)$.
Take partial derivatives over $T(L)$ in the remaining $\Phi_\sigma$.
They commute with $D$.
The constant-field assertion gives
\[
 \partial_{-\eta_j,0}y=Q_j(L),\quad
 \frac{dQ_j}{dL}=p_j,\qquad
 \partial_\sigma y=c_\sigma\in\C
 \quad(\sigma\in\mathfrak S_0\setminus\mathfrak U).
\]
Thus
\[
\begin{aligned}
 y_0&=y-\sum_jQ_j(L)G_{\eta_j,1}
       -\sum_{\sigma\in\mathfrak S_0\setminus\mathfrak U}
           c_\sigma\Phi_\sigma\in T(L),\\
 Dy_0&=-\sum_jQ_j(L)e^{-\eta_jz}/z
       -\sum_{\sigma\in\mathfrak S_0\setminus\mathfrak U}
           c_\sigma\phi_\sigma.
\end{aligned}
\]
Lemma~\ref{lem:reduction} gives $y_0\in R_1$.
Choose $j$ with $p_j\ne0$.
The coefficient of weight $-\eta_j$ satisfies
\[
 Du-\eta_ju=-Q_j(L)/z+g(z),\qquad
 u\in\C(z,L),\quad g\in\C(z).
\]
Here $Q_j$ is nonconstant.
Distinct nodes give distinct weights.
Lemma~\ref{lem:twisted-obstruction} excludes this equation.

\emph{2. Independence and constants at positive height.}
Assume \eqref{eq:appendix-tower-assertions} at every smaller height.
Choose $j$ with $h_j\ge2$.
Let $e_j$ be the $j$th unit vector.
Put $\widetilde h=h-e_j$.
For $K=\mathcal N_{\widetilde h}$, its obstruction gives
\[
 a=G_{\eta_j,h_j-1}/z\notin DK.
\]
Use coefficient $1$ at $j$ and $0$ elsewhere.
Apply Lemma~\ref{lem:appendix-primitive-extension} to $Y=G_{\eta_j,h_j}$.
Since $DY=a$, independence and constants follow at $h$.
The fixed meromorphic domain also gives the constant assertion directly.

\emph{3. The obstruction at positive height.}
Put $\mathcal J_h=\{j:h_j\ge2\}$ and $h^-_j=\max(1,h_j-1)$.
The independence just proved gives
\[
 \mathcal N_h=\mathcal N_{h^-}(Y_j:j\in\mathcal J_h),\qquad
 Y_j=G_{\eta_j,h_j},\quad DY_j\in\mathcal N_{h^-}.
\]
Suppose $Dy=\mathcal T_h(\mathbf p)$ with $\mathbf p\ne0$.
Take partials in the $Y_j$ for $j\in\mathcal J_h$, fixing $\mathcal N_{h^-}$.
Their commutation with $D$ and the constant assertion give
\[
 \partial_jy=Q_j(L),\quad \frac{dQ_j}{dL}=p_j
 \quad(j\in\mathcal J_h),\qquad
 y_0=y-\sum_{j\in\mathcal J_h}Q_j(L)Y_j\in\mathcal N_{h^-}.
\]
Differentiation gives
\begin{equation}\label{eq:appendix-tower-descent}
 Dy_0=\mathcal T_{h^-}(R),\qquad
 R_j=\begin{cases}
 p_j,&h_j=1,\\
 -Q_j,&h_j\ge2.
 \end{cases}
\end{equation}
A nonzero $p_j$ at height one survives.
Otherwise some $Q_j$ has nonzero derivative.
Thus $R\ne0$, while $H(h^-)<H(h)$.
The induction hypothesis excludes \eqref{eq:appendix-tower-descent}.

\emph{4. The listed solvable coordinates.}
Enlarge $\mathfrak S$ to contain their first-level supports.
Put $\mathcal P=T(\Phi_\sigma:\sigma\in\mathfrak S)$.
Choose a listed residue basis $B_1,\ldots,B_\rho$.
Lemma~\ref{lem:primitive} gives, for every other first-level coordinate,
\[
 F_a=\sum_i\kappa_{ai}B_i+r_a+c_a,\qquad
 \kappa_{ai},c_a\in\C,\quad r_a\in R_0.
\]
The higher $G_{\eta,k}$ are independent over $\mathcal P$.
Since $L\in\mathcal P$, so are $\Ein_k(\eta z)=L^k/k!-G_{\eta,k}$.
They remain independent over $T(B_1,\ldots,B_\rho)\subset\mathcal P$.
Eliminating the affine graphs leaves a polynomial ring in the $B_i$
and the $\tau$ distinct higher coordinates.
Hence the degree is $\rho+\tau$ and these graphs generate the ideal.


\subsection{The relation ideal}\label{app:ideal-details}
Retain the hypotheses of Lemma~\ref{lem:ideal}.
Put
\[
 \mathcal R=\C(z)[\Theta^{\pm1}],\qquad n=\rho+\tau,
 \qquad m=\#\mathcal C.
\]
Let $\mathcal I$ be the ideal of the stated generators.
The duplicate, constant, exponential, residue and linkage relations
are monic graphs. Successive substitution eliminates each dependent
variable. The fourth-coordinate shifts in \eqref{eq:core-quadric}
then give the following quotient.
\begin{equation}\label{eq:app-ideal-model}
 B_{\mathcal R}=
 \frac{\mathcal R[Z_1,\ldots,Z_n,(X_{\mathfrak c},Y_{\mathfrak c},
 S_{\mathfrak c},W_{\mathfrak c})_{\mathfrak c\in\mathcal C}]}
 {(X_{\mathfrak c}W_{\mathfrak c}-Y_{\mathfrak c}S_{\mathfrak c}
       +v_{\mathfrak c}:\mathfrak c\in\mathcal C)},
 \qquad v_{\mathfrak c}=u_{\mathfrak c}/z\in\mathcal R^\times.
\end{equation}
Here $W_{\mathfrak c}$ is the shifted fourth variable.
The variable $X_{\mathfrak c}$ represents $y_{\mathfrak c}$.

Lemma~\ref{lem:quadric} gives an $\mathcal R$-basis: use the $Z$-monomials
and, in each block, monomials avoiding $X_{\mathfrak c}W_{\mathfrak c}$.
Thus $B_{\mathcal R}$ is free over $\mathcal R$. It injects into its generic fibre.
That fibre is geometrically integral by Lemma~\ref{lem:quadric}
and \cite[Proposition~(4.6.5)(ii), p.~69]{EGAIV2}.
Hence $B_{\mathcal R}$ is a domain.

Put $M=\prod_{\mathfrak c}X_{\mathfrak c}$.
The same basis shows $M\ne0$. Solving each determinant gives
\begin{equation}\label{eq:app-ideal-chart}
 B_{\mathcal R}[M^{-1}]
 =\mathcal R[Z_1,\ldots,Z_n,
      (X_{\mathfrak c}^{\pm1},Y_{\mathfrak c},S_{\mathfrak c})_
        {\mathfrak c\in\mathcal C}],
 \qquad
 W_{\mathfrak c}=\frac{Y_{\mathfrak c}S_{\mathfrak c}-v_{\mathfrak c}}
 {X_{\mathfrak c}}.
\end{equation}
The earlier degree proof makes these selected functions independent over
$T$. Lemma~\ref{lem:exponential-log} identifies
$\operatorname{Frac}\mathcal R$ with $T$.
Evaluation on \eqref{eq:app-ideal-chart} is therefore injective.
Its restriction to $B_{\mathcal R}$ is injective too.
Consequently $\mathcal I$ is the full Laurent kernel and is prime.
Localization in $\mathcal R\setminus\{0\}$ gives the kernel over $T$.

Let $k_1\subset\C$ contain the generator data and residue constants.
The same substitutions and monic divisions work over
$\mathcal R_1=k_1(z)[\Theta^{\pm1}]$.
Their monomial basis maps to the identical basis over $\mathcal R$.
Thus $B_{\mathcal R_1}\hookrightarrow B_{\mathcal R}$.
Evaluation remains injective, proving descent of the full kernel.
Its fraction field is rational in the selected variables over
$k_1(z)(\Theta)$. Its transcendence degree over $k_1(z)$ is $d+n+3m$.

For ordinary polynomial variables, adjoin $V_h$ with
$\Theta_hV_h-1=0$. Replace negative powers by $V_h$ and eliminate
$\Theta,V$. The quotient is the Laurent quotient above.
Hence contraction gives the full coordinate kernel.
Equivalently, clear negative powers, saturate by
$\prod_h\Theta_h$, and eliminate.
Listed exponential variables require the corresponding saturation too.

\subsection{Specialization at nonzero algebraic points}
\label{app:specialization-details}
Assume the hypotheses of Theorem~\ref{thm:values}.
Thus the list is closed, normalized and saturated.
Its rates, kernels, scales and constants are algebraic.
Its Bessel orders are rational; its Kummer parameters are rational
and nonintegral. Scales are nonzero.
The primitives satisfy $(\mathrm{B2}^E)$, including $F(0)=0$.
First replace integer-order blocks by their order-zero blocks.
Put
\[
 R=\Qbar[z^{\pm1}][\Theta^{\pm1}],\qquad
 \mathcal R_\alpha=\Qbar[\Theta^{\pm1}],\qquad
 n_{\mathrm{val}}=d+\rho^E+\tau+3m,\quad m=\#\mathcal C.
\]

The residue vectors have algebraic entries. Hence their basis
coefficients $\kappa_{ji}$ are algebraic.
For $\mu\in\Qbar$, put $D_\mu g=g'+\mu g$.
On $\Qbar[z^{\pm1}]$,
\begin{equation}\label{eq:app-laurent-image}
 \operatorname{im}D_\mu
 =\ker\bigl(q\mapsto\Res_0 e^{\mu z}q(z)\,dz\bigr).
\end{equation}
The residue vanishes on $D_\mu g$ since
$e^{\mu z}D_\mu g=(e^{\mu z}g)'$.
For $\mu=0$, differentiate Laurent monomials.
For $\mu\ne0$, $D_\mu$ is invertible on polynomials.
Its negative-power reduction is
$z^{-j-1}\equiv(\mu/j)z^{-j}$ for $j\ge1$.
The remaining class is $z^{-1}$, with residue $1$.
This proves \eqref{eq:app-laurent-image}.
Apply it to each residue-zero first-level relation.
It gives $r_j\in R$ in \eqref{eq:residue-relations}.
Zero-origin normalization gives
\[
 c_j=-\lim_{z\to0}r_j(z,e^{\lambda_1z},\ldots,e^{\lambda_dz})
     \in\Qbar.
\]
The limit exists because the corresponding primitive difference is entire.
It is a finite Laurent--exponential constant coefficient.

All other graph coefficients lie in $R$.
For the shift $(f_\nu,f'_\nu)\mapsto(f_{\nu+1},f'_{\nu+1})$ and the shift
$(f_{-\nu},f'_{-\nu})\mapsto(f_{-\nu-1},f'_{-\nu-1})$, respectively,
the solution-coordinate determinants are
\[
 -\frac{4(\nu+1)^2}{c^2z^2},\qquad
 -\frac{c^2z^2}{4\nu^2}.
\]
Their product is $(\nu+1)^2/\nu^2$.
The normalized noninteger shifts have $\nu(\nu+1)\ne0$.
Kummer contiguity has determinant $1$ on each diagonal block.
The exponential links have unit determinants in $R$.
Integer-order reduction has determinant $(-1)^n$ on each diagonal block
by Lemma~\ref{lem:integer}.
Scale and order signs use identities or permutations.
Thus every required map and inverse specializes at $\alpha\ne0$.

Appendix~\ref{app:arithmetic} supplies the $E$-function bounds.
Use the vector in \eqref{eq:value-system} after order-zero reduction.
Its primitive equations have Laurent-polynomial coefficients.
Its tower equations are \eqref{eq:ein-tower}.
The block equations are
\[
\begin{gathered}
 I''=c^2I-I'/z,\qquad
 S''=c^2S-S'/z+2I'/z,\\
 F''=(s-z^{-1})F'+asF/z,\\
 \Sigma''=(s-z^{-1})\Sigma'+as\Sigma/z+(2F'-sF)/z,\\
 f''_{\nu,c}=c^2f_{\nu,c}-(2\nu+1)f'_{\nu,c}/z.
\end{gathered}
\]
Consequently $f$ satisfies a first-order system with coefficient matrix
in $\Qbar[z,z^{-1}]$.
Every $\alpha\in\Qbar^\times$ is ordinary.
Each component of $f$ has algebraic Taylor coefficients. Any polynomial
relation over $\C(z)$ therefore descends to $\Qbar(z)$ by solving its
coefficient equations over $\Qbar$; the two transcendence degrees agree.
Lemma~\ref{lem:ideal} and the generic degree formula give
\begin{equation}\label{eq:app-specialized-degree}
 \trdeg_{\Qbar}\Qbar(f(\alpha))
 =\trdeg_{\Qbar(z)}\Qbar(z)(f)=n_{\mathrm{val}}
\end{equation}
by \cite[Theorem~1.1, p.~370]{Beukers2006}.
The extra system exponentials are Laurent monomials in the basis
exponentials. Thus they leave both fields unchanged.

Specialize the graph eliminations at $z=\alpha$.
They give \eqref{eq:app-ideal-model} over $\mathcal R_\alpha$, with
$n=\rho^E+\tau$ and
\[
 v_{\mathfrak c}\in
 \{\alpha^{-1},\Theta^{v_s}\alpha^{-1},2\nu\alpha^{-1}\}
 \subset \mathcal R_\alpha^\times.
\]
Write $B_\alpha$ for this quotient.
The free-basis argument above makes $B_\alpha$ a domain.
Its generators vanish on the values.

Let $s_0$ count the blocks with $X_{\mathfrak c}(\alpha)=0$.
In each such block,
\[
 Y_{\mathfrak c}(\alpha)S_{\mathfrak c}(\alpha)
   =v_{\mathfrak c}(e^{\lambda_1\alpha},\ldots,e^{\lambda_d\alpha})
   \ne0.
\]
Over the exponential-value field, this block needs two generators:
$Y_{\mathfrak c}(\alpha),W_{\mathfrak c}(\alpha)$.
Every other block needs three, by its determinant.
All values are therefore generated by at most $n_{\mathrm{val}}-s_0$ numbers.
Equation~\eqref{eq:app-specialized-degree} forces $s_0=0$.
It then makes the $n_{\mathrm{val}}$ selected values, including basis exponentials,
algebraically independent.

Evaluation now extends to $B_\alpha[M^{-1}]$.
By \eqref{eq:app-ideal-chart}, this is a Laurent polynomial ring
in those selected variables. Its evaluation is injective.
Since $B_\alpha\hookrightarrow B_\alpha[M^{-1}]$, evaluation on
$B_\alpha$ is injective. Hence $I_\alpha=\mathfrak P_\alpha$.
In particular $y_{\mathfrak c}(\alpha)\ne0$.
Nonvanishing was established before localization.
The system denominators are powers of $z$; all coordinate maps have
unit determinants in $R$. These facts account for every excluded point.

For several points put $\widetilde q_i(z)=q_i(\alpha_i z)$.
Primitive kernels become $\alpha_iq_{i,\mu}(\alpha_i z)$.
Residues retain their values under
$(\mu,p)\mapsto(\alpha_i\mu,p/\alpha_i)$.
Tower pairs become $(\alpha_i\eta_i,k)$.
Derivative coordinates acquire the nonzero factors $\alpha_i$.
Close and normalize the combined list; saturate its rates.
Test residues, tower pairs and linkage classes on this list.
The one-point result at $1$ gives the several-point assertions.


\label{page:appendix-end}
\clearpage
\section*{Supplementary material}
\setcounter{section}{0}
\renewcommand{\thesection}{S\arabic{section}}
\renewcommand{\theHsection}{supplement.\arabic{section}}
This supplementary material contains the mixed example, exact input model,
measure comparison, normalized-function example, source comparison and
coefficient-cancellation details.
\section{A mixed example}\label{sup:mixed-example}

Let $\Gamma=\Z$. Take $e^z$, the tower levels
$\Ein_k(z)$ ($1\le k\le3$), $\Ein_1(-2z)$, and the
origin-normalized primitive
\[
 (F^E)'=(e^z-1-z)/z^2,\qquad F^E(0)=0.
\]
For the second-order part, use nine blocks. The representative Bessel parameters $(c,\nu)$ are
$(1,0)$, $(i,0)$ and $(1,1/3)$. The representative Kummer parameters
$(s,a)$ are $(-1,2/3)$ and $(2,1/3)$. Add $\operatorname{Bess}(1,1)$.
Complete the list with the three Kummer parameters $(2,1/2)$, $(1,1/3)$ and $(1,4/3)$.
Retain the four slots of the integer-order block.
The residue vectors are $e_{(0,0)}-e_{(-1,0)}$,
$e_{(0,0)}-e_{(2,0)}$, and $e_{(1,0)}-e_{(0,0)}$.
They are independent and have total sum zero. Thus $e_{(0,0)}$ is outside
their span. The logarithm is unlisted. Thus $d=1$, $\rho=3$, $\tau=2$ and
\begin{equation}\label{eq:mixed-example-degree}
 \trdeg_{\C(z)}K_{\rm coord}=1+3+2+3\cdot5=21.
\end{equation}
The list has $42$ coordinate slots. Adjoining $L$ raises the residue rank to
four and the degree to $22$, in agreement with main-text Theorem~\ref{thm:main}.
The ancillary script \texttt{anc/n1\_exact.py} and output
\texttt{anc/n1\_exact.out} verify both presentations and their saturation
in the formal coordinate ring over $\Q(z)$. See \texttt{anc/N1-README.md}.

\section{The input model and computation}\label{sup:input-model}
The complete definition, computability proposition and proof follow
the overview below.
A presented $k_0\subset\C$ has effective arithmetic and equality with a
specified embedding. In model (M-alg), a minimal polynomial and isolating
disc specify a number field. Model (M-gen) supplies pole factors and oracle (O1), returning
a $\Z$-basis of integer relations in any finite input tuple.
Rates, nodes, scales and Kummer parameters lie in $k_0$; Bessel orders
are reduced rationals and kernels are polynomial pairs. Put $k_1=k_0(\text{poles})$. A base-point input specifies
$F=\sum_{\sigma=(\mu,p)}e^{\mu p}\operatorname{res}(F)_\sigma\Phi_\sigma+r_F+\gamma_F$
with $\gamma_F\in k_0$. At weight zero the polynomial part of $r_F$ has zero
constant. Origin-normalized inputs instead determine the constant by the
limit at zero. Accordingly, a fundamental residue relation has constant
$c_j=\gamma_j-\sum_i\kappa_{ji}\gamma_i$ on the base-point side and
$c_j=-\lim_0r_j(z,e^{\lambda_1z},\ldots,e^{\lambda_dz})$ on the origin side.
Compatibility requires the two residue spans to intersect trivially after
supplied exact normalization comparisons. These encode the base point. Thus $z_*$ is not used directly. Their correctness is assumed because
analytic evaluation is excluded. For compatible inputs, integer/field
linear algebra, rational reduction and Gr\"obner elimination compute
$d,\rho,\tau,\#\mathcal C$, degrees, generators and all sublist ideals
over $k_1(z)[\Theta^{\pm1}]$. Incompatible inputs return no ideal. Approximations alone are excluded. In model (M-alg), integer linear
algebra supplies the integer relations. Thus no oracle is required.
Values require (M-alg), $\alpha\in k_0^\times$ and main-text Theorem~\ref{thm:values}'s
hypotheses. Retain distinct coordinates, emit duplicate graphs and $X_c-c$,
include the half-integer exponential graphs with rates $\pm c$, then compute
the saturated lattice, residues and normalization constants, and eliminate
inverse variables using $\Theta_h\Theta_h^{-1}-1$.
Contraction gives each sublist ideal and degree. All operations terminate
in the specified field without analytic evaluation. Implementation details
are supplementary. The value coefficient bounds are printed in
main-text Appendix~\ref{app:arithmetic}.
\subsection{Exact input definition and computation}\label{sup:input-algorithm}
\label{subsec:input}
\begin{definition}\label{def:input}
A presented field $k_0\subset\C$ has finite representations for its elements.
Arithmetic and equality are computable. Its embedding is specified.
Model (M-alg) supplies a number field $\Q(\theta)$.
A minimal polynomial and isolating disc specify $\theta$.
In model (M-gen), oracle (O1) takes any finite input tuple.
It returns a $\Z$-basis of its integer relations.
For number fields this oracle is supplied by integer linear algebra.
To test $a\in\Z$, ask whether there is a relation $ma+n=0$ with
$m\ne0$ and $m\mid n$.

The rates, nodes, scales and Kummer parameters belong to $k_0$.
Bessel orders are supplied as reduced rationals.
Rational kernels are supplied as polynomial pairs.
Model (M-gen) also supplies their linear denominator factors.
In (M-alg), we compute a splitting field.
Set $k_1=k_0(\text{poles})$. The value statements require (M-alg) and
$\alpha\in k_0^\times$. The supplied constants $\gamma_F$ encode
the chosen base-point normalization. Thus $z_*$ does not enter the computation.
A base-point primitive is specified by $\gamma_F\in k_0$:
\begin{equation}\label{eq:primitive-normalization}
F=\sum_{\sigma=(\mu,p)}e^{\mu p}\operatorname{res}(F)_\sigma\Phi_\sigma+r_F+\gamma_F.
\end{equation}
At nonzero weights, rational reduction uniquely determines $r_F$.
At weight zero, its polynomial part must have zero constant.
Thus $\gamma_L=0$. Coordinates of type (B2$^E$) and towers instead use
origin normalization. Let $\mathcal V_{{\rm res},0}$ be the origin-normalized first-level residue span.
Write $\mathcal V_{{\rm res},*}$ for the base-point-normalized first-level residue span.
The list is normalization-compatible when
\begin{equation}\label{eq:compatibility}
\mathcal V_{{\rm res},0}\cap\mathcal V_{{\rm res},*}=0,
\end{equation}
or after moving a specified origin subset to the base-point side using
supplied $\gamma_G\in k_0$ satisfying \eqref{eq:primitive-normalization}.
A supplied comparison is assumed correct; its constant is exact input.
The algorithm does not evaluate the analytic primitive at the base point.
Choose residue bases separately on the two
normalization sides. Each fundamental coordinate-basis relation uses one normalization.
An arbitrary relation decomposes into the two types.
Compatibility holds if all coordinates except possibly $L$ are origin-normalized.
Their residue vectors have $\sum_\mu v_{(\mu,0)}=0$.
For $L$, $\operatorname{res}(L)=e_{(0,0)}$.
It also holds if every coordinate uses the base point.
\end{definition}

\begin{proposition}\label{prop:algorithm}
In both input models, the quantities $d,\rho,\tau,\#\mathcal C$ and the
corresponding degrees are computable. For compatible lists, the generators
and every sublist ideal are computable over $k_1(z)[\Theta^{\pm1}]$.
Model (M-alg) requires no oracle, and (M-gen) uses (O1). For incompatible
lists, the algorithm reports undetermined mixed constants and returns no
ideal. Inputs supplied only by complex approximations are excluded.
For values, the hypotheses of main-text Theorem~\ref{thm:values} are required.
\end{proposition}
\begin{proof}

The algorithm first rejects inadmissible specifications. It then closes
and normalizes the list, retaining one representative per coordinate.
It emits $X_q-X_{\operatorname{rep}(q)}$ for every duplicate slot of every type.
It emits $X_c-c$ for each listed constant $c\in k_1$.
Only distinct higher levels are free.
Half-integer Bessel coordinates retain main-text Section~\ref{sec:main-statement}'s exponential graphs.
They also retain its derivative/order-shift graphs.
The saturated lattice contains their rates $\pm c$.
They contribute no class to $\mathcal C$.
The algorithm preserves these graphs through closure and specialization.
It also preserves them through saturation and contraction.
It composes the linkage maps. Oracle (O1), followed by
Hermite normal form, gives the minimal saturated lattice. Partial
fractions and finite Taylor expansions determine the residues. Linear algebra then gives $\rho$, $\kappa_{ji}$ and the test for
$e_{(0,0)}\in\mathcal V_{\rm res}$; $\tau=\sum_\eta(h_\eta-1)$ is the number of distinct retained levels $(\eta,k)$ with $k\ge2$. Rational reduction with the specified
normalization gives
\begin{equation}\label{eq:relation-constants}
c_j=\begin{cases}\gamma_j-\sum_i\kappa_{ji}\gamma_i,&\text{base-point},\\
-\lim_0r_j(z,e^{\lambda_1z},\ldots,e^{\lambda_dz}),&\text{origin}.
\end{cases}
\end{equation}
Any supplied normalization comparisons are applied before this step.
Next adjoin inverse variables with relations
$\Theta_h\Theta_h^{-1}-1$, saturate, and eliminate. Contraction gives
the sublist ideals and the degrees of the original coordinate lists.
For values, specialize only $z=\alpha$.
The procedure uses field arithmetic and equality, with (O1).
It also uses Hermite normal form and Buchberger elimination.
For compatible lists, these operations obtain all required constants.
No analytic evaluation is used.
\end{proof}

The arithmetic coefficient bounds are printed in main-text Appendix~\ref{app:arithmetic}.


\section{Comparison of linear independence estimates}\label{sup:measure-comparison}
Main-text Theorem~\ref{thm:measure} has an explicit threshold $H_0$.
It is too large for practical small-height tests.
Galochkin's measure \cite[theorem, p.~241]{Galochkin1968} applies to the
algebraically independent pair $(e^{-z},\Ein(z))$. Multiplication of our
form by $e^{-1}$ gives the polynomial of degree at most two $b+ax+cx^2+dxy$ at
$(x,y)=(e^{-1},\Ein(1))$. This gives a power-type bound with unevaluated
constant. Bertrand--Beukers \cite[Theorem~1, p.~182]{BertrandBeukers1985}
give the multiplicity estimate for this four-function system;
Bertrand--Chirskii--Yebbou \cite[Theorem~1.2, p.~245]{BertrandChirskiiYebbou2004}
make its Shidlovskii constant effectively computable.
These are inputs to effective measures.
They give no numerical constants for this quadruple.
Shidlovskii's linear independence measure
\cite[Chapter~11, \S2, Theorem~1, (32), p.~357]{Shidlovskii1989}
gives the following bound for the rational system $(1,e^z,e^{-z},\Ein(z))$
at the ordinary point $1$: for every $\varepsilon>0$ there is
$c_\varepsilon>0$, depending only on $\varepsilon$ and this fixed system
and point, such that every $(a,b,c,d)\in\Z^4\setminus\{0\}$ satisfies
$\lvert a+be+ce^{-1}+d\Ein(1)\rvert>c_\varepsilon H^{-3-\varepsilon}$,
where $H=\max\{\lvert a\rvert,\lvert b\rvert,\lvert c\rvert,\lvert d\rvert\}\ge1$.
Functional linear independence is required.
The proof is on pp.~357--359.
The book does not evaluate the positive constant.
Fischler--Rivoal state the rational case as Theorem~A.
Their extension is Theorem~1 \cite[pp.~2--3]{FischlerRivoal2024}.
Page~3 discusses effectiveness in principle.
It does not compute the constants.
Their Theorem~1 assumes a nonzero linear form.
For $(a,b,c,d)\in\Z^4\setminus\{0\}$, Theorem~\ref{thm:values}
implies $a+be+ce^{-1}+d\Ein(1)\ne0$.
Main-text Theorem~\ref{thm:measure} evaluates the constants.
Its exponent $-60$ is weaker.
The threshold $H_0=N^N$ is astronomically large.
The displayed bound applies for $H\ge H_0$. Below $H_0$, scaling the coefficient vector
gives only $|\Lambda|\ge H/(12(2H_0)^{61})$.
Main-text Appendix~\ref{sec:measure} evaluates the $\nu=\mu=4$ polynomial-and-derivative
construction of Adamczewski--Faverjon.
A direct cyclic-vector Wronskian bound replaces a general zero estimate.
This section records the comparison with Galochkin's measure
and the multiplicity estimates of Bertrand--Beukers and
Bertrand--Chirskii--Yebbou.
For classical background, Fel'dman--Nesterenko's monograph
\cite{FeldmanNesterenko1998} is a further reference. Beukers points to its
Theorem~5.23 for the qualitative Siegel--Shidlovskii theorem
\cite[p.~369]{Beukers2006}.
This comparison attributes no further exponent or evaluated constant to the monograph.

\section{A first example}\label{sup:first-example}
For $c\ne0$, the normalized maps and determinant in the introductory
example are
\[
\begin{gathered}
 f_{1/2}(2cz)=e^{cz}I_c(z),\qquad
 S_{1/2}(2cz)=e^{cz}S_c(z),\\
 z(I_cW_c-V_cS_c)-I_c^2+1=0.
\end{gathered}
\]
Their derivatives identify $\operatorname{Bess}(c,0)$ and
$\operatorname{Kum}(2c,1/2)$ as one class. The following formulas use $c=1$.
Put $\mathsf H_0=0$ and $\mathsf H_n=\sum_{j=1}^n1/j$. Consider the four entire functions
\begin{equation}\label{eq:supp-intro-series}
\begin{aligned}
 I(z)&=\sum_{n\ge0}\frac{(z/2)^{2n}}{(n!)^2},&
 \mathsf S(z)&=\sum_{n\ge0}\frac{\mathsf H_n(z/2)^{2n}}{(n!)^2},\\
 f_a(u)&=\sum_{n\ge0}\frac{(a)_nu^n}{(n!)^2},&
 S_a(u)&=\sum_{n\ge0}\frac{(a)_nh_n(a)u^n}{(n!)^2},
\end{aligned}
\end{equation}
where $(a)_0=1$, $(a)_n=a(a+1)\cdots(a+n-1)$, and
$h_n(a)=2\mathsf H_n-\sum_{j=0}^{n-1}(j+a)^{-1}$. The parameter $a$ is
nonintegral. The half-Kummer identities are
\begin{equation}\label{eq:supp-intro-half}
 f_{1/2}(2z)=e^zI(z),\qquad S_{1/2}(2z)=e^z\mathsf S(z).
\end{equation}
These classical identities \cite[10.39.5,13.2.9,13.6.10]{DLMF}
are proved by operators in main-text Section~\ref{sec:linkage}.

These identities also explain cancellation of Euler's constant. Suppose $x>0$.
The standard Bessel function $K_0$ and Kummer function $U$ satisfy
\begin{equation}\label{eq:supp-intro-connection}
\begin{aligned}
 \frac{K_0(x)}{I(x)}
   &=\frac{\mathsf S(x)}{I(x)}-\gamma-\log(x/2),\\
 \frac{\sqrt\pi\,U(1/2,1,2x)}{f_{1/2}(2x)}
   &=\frac{S_{1/2}(2x)}{f_{1/2}(2x)}-\gamma-\log(x/2).
\end{aligned}
\end{equation}
With the normalizations of \cite[10.31.2, 13.2.9, 13.6.10]{DLMF},
the denominators are positive. Subtraction cancels both connection terms:
\begin{equation}\label{eq:supp-intro-gamma-free}
 \frac{\sqrt\pi\,U(1/2,1,2x)}{f_{1/2}(2x)}
 -\frac{K_0(x)}{I(x)}
 =\frac{S_{1/2}(2x)}{f_{1/2}(2x)}-\frac{\mathsf S(x)}{I(x)}=0.
\end{equation}
Put $D=d/dz$. The determinant relation
\begin{equation}\label{eq:supp-intro-determinant}
 z\{I\,D\mathsf S-(DI)\mathsf S\}-I^2+1=0
\end{equation}
leaves three independent Bessel coordinates. Put
\[ (Q_1,Q_2,Q_3)=(f_{1/2}(2z),D[f_{1/2}(2z)],S_{1/2}(2z)). \]
Thus $\C(z)(I,DI,\mathsf S,Q_1,Q_2,Q_3)=\C(z)(e^z,I,DI,\mathsf S)$ by
\eqref{eq:supp-intro-half} and its derivative. Its degree is four over
$\C(z)$ and three over $T=\C(z)(e^z)$.
The blocks form one linkage class.

\section{Classical single-family results}\label{sup:classical-comparison}
For one scale, the resonant Bessel triple is classical. Mahler proved
independence of values of $J_0,J'_0$ and order derivatives of normalized
$J_\lambda$ at $\lambda=0$ \cite[Theorem~2, p.~167]{Mahler1968}. His
Theorem~3, pp.~169--170, rewrites these values with connection constants.
In particular, he proves that $\pi Y_0(2)/(2J_0(2))-\gamma$ is
transcendental \cite[Theorem~3 and its consequence, pp.~169--170]{Mahler1968}.
His $K_\lambda$ denotes normalized ordinary Bessel $J$, not modified $K$.
Belogrivov treated each
nonnegative integer order at nonzero algebraic points and the determinant
relation $z(u'_1u_0-u_1u'_0)+u_0^2=1$
\cite[Theorem~4, Lemmas~4--5, (36)]{Belogrivov1973}.
His coordinates are $u_{0,\lambda}(\zeta)=K_\lambda(\zeta)$ and
$u_{1,\lambda}(\zeta)=\partial_\lambda K_\lambda(\zeta)$, with the
same normalized ordinary Bessel convention. Thus
$u_{0,0}(icz)=I_c(z)$ and $u_{1,0}(icz)=-S_c(z)$.
His primes are derivatives in $\zeta$; substitution gives
$z(I_cW_c-V_cS_c)-I_c^2+1=0$.
His Kummer-pair theorems treat values of ${}_1F_1(\mu;\nu;z)$ and its
derivative at several points, without the logarithmic companion
\cite[Theorems~1--2]{Belogrivov1971}.
The one-point corollary over the full $\mathcal E$ is known from
Shidlovskii's simultaneous exponential-integral theorem
\cite[Chapter~7, \S1, Theorem~1, pp.~232--233]{Shidlovskii1989}
by the finite-rate reduction below.
Mahler's final integral family was announced without proof
\cite[\S23, pp.~172--173]{Mahler1968}, as Rivoal notes
\cite[\S7]{Rivoal2012}.
Put $G_0(z)=\int_0^\infty e^{-t}(t+z)^{-1}\,dt$ for
$z\in\C\setminus(-\infty,0]$, and use the principal branch of $\log z$
\cite[p.~1 of the author manuscript]{Rivoal2012}.
Rivoal proves his Theorem~2(ii) for every algebraic
$z\notin(-\infty,0]$ \cite[Theorem~2(ii), \S\S5 and~7]{Rivoal2012}:
$\trdeg_{\Q}\Q(\gamma+\log z,e^z,G_0(z))\ge2$.
At $z=1$, $G_0(1)=\delta=eE_1(1)$, so this gives
$\trdeg_{\Q}\Q(e,\gamma,\delta)\ge2$.
The full $\mathcal E$-base conclusion is attributed to Shidlovskii's
simultaneous theorem, without adding an all-rate quantifier to Rivoal's statement.
Main-text Theorem~\ref{thm:connection} systematically assembles
multi-node mixed families over the selected regular-coordinate field $B$.
The source comparison for that scope remains partial.

\subsection*{Finite-rate reduction for the known one-point consequence}
Shidlovskii defines $\omega_0(z)=e^z$ and
$\omega_k(z)=1+\sum_{n\ge1}z^n/(n!n^k)$ for $k\ge1$.
Thus $\omega_1(-x)=1-\Ein(x)$.
For each integer $m\ge1$ and each finite set of algebraic nodes
$\xi_1,\ldots,\xi_r$ linearly independent over $\Q$, his theorem gives
simultaneous algebraic independence of $\omega_k(\xi_i)$,
$0\le k\le m$, $1\le i\le r$. The reduction below uses $m=1$.
Fix algebraic $x\ne0$. A polynomial relation for $\Ein(x)$ over
$\mathcal E$ has finitely many coefficients, each algebraic over a
field generated by finitely many $e^{\beta_j}$ with algebraic $\beta_j$.
Extend $\eta_1=-x$ to a $\Q$-basis $\eta_1,\ldots,\eta_r$ of the
span of $-x$ and these rates.
Clearing rational denominators makes every $e^{\beta_j}$ algebraic
over $K=\Qbar(e^{\eta_1},\ldots,e^{\eta_r})$.
The theorem with $m=1$ makes $\omega_1(-x)$ transcendental over $K$.
It remains transcendental over any algebraic extension containing the
finitely many coefficients, excluding the proposed relation.
At $x=1$, $\Ein(1)=\gamma+E_1(1)$ and $e\in\mathcal E^\times$,
giving the equivalent alternatives with $E_1(1)$ and $\delta=eE_1(1)$.
This reduction retains $-x$ as a basis member. It does not retain
arbitrary merely distinct nodes in a $\Q$-basis or adjoin regular
Bessel/Kummer pairs; the mixed conclusion over $B$ uses
main-text Theorem~\ref{thm:values}.

Powers treats odd and even multisections of iterated exponential integrals.
He proves independence over the field generated by their level-zero functions.
At $1$, he proves independence over the corresponding value field \cite[p.~13]{Powers2026}.
Belogrivov also treated one-node shifted-denominator hypergeometric
families and argument derivatives \cite{Belogrivov1970}.

\subsection*{Homogeneous and integration antecedents}
Ostrowski makes primitive dependence explicit for rational-exponential
derivatives \cite[pp.~315--316]{Ostrowski1946}; Kolchin gives the general
differential-algebraic setting \cite{Kolchin1973}.
Beukers, Brownawell and Heckman's 1988 paper
\cite{BBH1988} is an antecedent for the differential-Galois method.
BBH's homogeneous framework supplies its normality characterization
\cite[Theorem~2.2, p.~288]{BBH1988}.
Its common-power descent assumes irreducible Siegel-normal systems
\cite[Proposition~2.4, pp.~291--292]{BBH1988}.
Gorelov invokes these results in his homogeneous reduction
\cite[Lemma~13, p.~7]{Gorelov2021}.
He identifies BBH's omitted Bessel--Kummer linkage
\cite[Example~1, p.~3; Lemma~14, p.~7]{Gorelov2021}.
The second BBH result cited there is Proposition~2.4, not Theorem~2.4.
His criteria characterize independence of fundamental-matrix entries with exponentials
and of the values $F^{(i)}(\alpha)$ \cite[Theorems~1--2, pp.~4--5]{Gorelov2021}.
Appendix~\ref{app:gorelov} checks the hypotheses for the homogeneous rank-two inputs used here.
BBH Theorem~2.2(iv) requires an extension by scalars of $\SL_2$.
This is the rank-two group condition supplied by main-text Lemma~\ref{lem:groups}.
It includes the homogeneous logarithmic solution $J=S-I L$ or $J=\Sigma-FL$.
No normality assertion applies to the enlarged extension with $L$ and the E-parts separately.

Gorelov's inhomogeneous generalized hypergeometric equation has scalar forcing
$\prod_{j=1}^q(\lambda_j-1)$ \cite[p.~663]{Gorelov2016}.
His inhomogeneous hypergeometric criteria require homogeneous-group and parameter hypotheses
\cite[Theorems~1--3, pp.~664--665]{Gorelov2016}.
For $q>l$, his function criteria work over homogeneous fundamental-matrix fields.
They require a homogeneous group containing $\SL_q$
\cite[Theorems~1--2, p.~664]{Gorelov2016}.
Dependence occurs when a denominator parameter belongs to $\mathbb N$;
for a collection, equivalent parameter vectors at equal scales also cause dependence.
His value theorem requires rational parameters, odd $q\ge3$, irreducibility exclusions and nonzero algebraic arguments.
It also imposes four explicit collision conditions
\cite[Theorem~3, pp.~664--665]{Gorelov2016}.
These hypotheses are not conditions for our normalized logarithmic E-parts.
In our resonant encoding, $\lambda=(1,1)$ makes that scalar forcing vanish.
The E-parts instead have forcing $2I'$ or $2F'-sF$.
Their name alone does not identify the two inhomogeneous problems.

Gorelov's primitive dependence criterion is a constant-linear-combination criterion
\cite[Lemma~6, p.~670]{Gorelov2016}.
His integration result adjoins primitives of $z^\delta e^{\omega z}$
over the specified homogeneous fields
\cite[Lemma~7, pp.~670--672]{Gorelov2016}.
Here $\delta$ is his integration exponent, distinct from the Gompertz constant above.
The nonzero rates are distinct, $\delta\notin\Z_{\ge0}$,
and same-rate exponents do not differ by integers.
Independence of his fundamental-matrix, exponential and power functions is assumed.
The case $\delta=-1$ includes first-level exponential integrals.
Lemma~8 concerns integration involving homogeneous solution cofactors.
These are integration antecedents.
They do not state the full finite-pole residue algorithm or its affine normalization constants.

BBH already discusses successive exponential integrals as a Liouville extension
\cite[p.~281]{BBH1988}.
Its successive series are $\sum_{n\ge1}z^n/(n^j n!)$ through finite depth.
They arise from successive inhomogeneous first-order equations, outside Siegel normality.
This anticipates tower integration, without giving the present multi-node tower ideal or its degree.

\subsection*{Gorelov's normalized integral family}
For $\lambda\in\C\setminus\Z_{<0}$, write
$\phi_\lambda(z)={}_1F_1(1;\lambda+1;z)=\sum_{n\ge0}z^n/(\lambda+1)_n$,
so $\phi_\lambda(0)=1$ \cite[p.~384]{Gorelov2025}.
Gorelov's later paper treats the specific family
$V_{\lambda,\nu,\alpha}(z)=\allowbreak\nu z^{-\nu}e^{\alpha z}\allowbreak\int_0^z t^{\nu-1}e^{-\alpha t}\phi_\lambda(t)\,dt$
\cite[pp.~383--384]{Gorelov2025}.
Here $\alpha\in\C$ and $\lambda,\nu\in\C\setminus\Z_{<0}$.
Explicitly, if $e^{-\alpha z}\phi_\lambda(z)=\sum_{n\ge0}a_nz^n$, then
$V_{\lambda,\nu,\alpha}(z)=e^{\alpha z}
\bigl(a_0+\sum_{n\ge1}\nu a_nz^n/(\nu+n)\bigr)$.
This is Shidlovskii's normalized primitive for $\nu\notin\Z_{<0}$
\cite[p.~384]{Gorelov2025}; in particular, $V_{\lambda,0,\alpha}=e^{\alpha z}$.
The integral need not converge at zero for every allowed complex parameter.
This family has $V_{\lambda,\nu,\alpha}(0)=1$.
It satisfies a second-order inhomogeneous equation and
$V'=(\alpha-\nu/z)V+(\nu/z)\phi_\lambda$
\cite[p.~384; (17), p.~385]{Gorelov2025}.
Theorem~1 classifies its Laurent-polynomial representations in exponentials and $\phi$ functions
for the stated complex parameters. For members of this family that are E-functions,
Theorem~2 classifies algebraic dependence of $V$ and $V'$ over $\C(z)$
\cite[pp.~384--385]{Gorelov2025}.
Lemma~3 gives a second-order inhomogeneous characterization of E-functions
\cite[p.~386]{Gorelov2025}.
Its value discussion concerns that family
\cite[Conclusion~1, p.~390]{Gorelov2025}.
Collections $V_{\lambda,\nu,\alpha}(\xi z)$ with varying parameters and scales are explicitly deferred
\cite[\S4, item~2 (Conclusion~2), p.~390]{Gorelov2025}.

These works supply homogeneous, integration and value antecedents.
They do not state the joint multi-scale normalized ideal with shared logarithmic E-parts,
finite-pole residue relations, higher Ein towers and fixed additive constants.
They also do not give the present actual connection-constant cancellation algebra.
No global novelty conclusion follows from this comparison.

\subsection*{Earlier value inputs}
Earlier value theorems have narrower inputs.
BBH treats rational admissible homogeneous hypergeometric families
\cite[Corollary~4.6, pp.~300--301]{BBH1988}.
Gorelov imposes rational-parameter and collision conditions
\cite[Theorem~3, pp.~664--665]{Gorelov2016}.
His later result concerns one integral family, with Theorem~2 restricted to E-functions
\cite[Theorem~2, p.~385; Conclusion~1, p.~390]{Gorelov2025}.
These statements do not supply the present specialized normalized ideal.

\section{Examples and limits of coefficient cancellation}\label{sup:cancellation-remarks}
The first two examples illustrate the general algebraic distinction.
For example, over algebraically closed $B$, take a transcendental $x$
and $X=x$, $d=x/2$. Then $V=x/2$ and $B[V]\cap B[X]=B[x]$, although
$\mathcal A_d=B$. If instead $d=2x$, then $P(Z)=Z^2$ gives
$P(V)=P(X)$ without being translation invariant. Thus independence of
$X$ over $B$ leaves mixed relations between $X$ and $\mathbf c$ undetermined,
including the arithmetic nature of $\gamma$.

For connection values, use main-text Theorem~\ref{thm:connection}'s notation.
Linked entries must first be reduced. Contiguity and the determinant give
\[
 \frac{\Gamma(a+1)U(a+1,1,t)}{f_{a+1}(t)}
 -\frac{\Gamma(a)U(a,1,t)}{f_a(t)}
 =-\frac{e^t}{a f_a(t)f_{a+1}(t)}\in B,
\]
and the half-Kummer entry at $(2u,1/2)$ equals $K_0(u)/I_0(u)$.
Higher Ein levels are excluded: their natural tail already satisfies
$\int_x^\infty e^{-v}\log(v/x)\,dv/v
 =\tfrac12(\log x+\gamma)^2+\pi^2/12-\Ein_2(x)$.
No reduction of this connection to $B+B\gamma$ is assumed.

Fischler--Rivoal's Gamma-polynomial connection mechanism
\cite[Theorem~2(iii) and the subsequent discussion on p.~35]{FischlerRivoal2016}
gives asymptotic coefficients at infinity.
They have the form $\Gamma(\alpha^*)P(\gamma)$ over their ring of $G$-values.
Here $\alpha^*=1$ for nonpositive integral $\alpha$.
Gauss's digamma formula and recurrence are classical.
Bundschuh~\cite[Corollary~1]{Bundschuh1979} proved transcendence of
$\psi(a)+\gamma$ for rational nonintegral $a$. See also
Murty--Saradha~\cite[Theorems~1, 3--4]{MurtySaradha2007}.
These results do not imply independence over the present $B$.

\subsection*{Detailed coefficient-cancellation argument}
Main-text Proposition~\ref{prop:coefficient-cancellation} is an algebraic
statement about the actual column $d$, independent of its special-function
origin. Write $d=d_0+M\xi$ with $d_0\in B^N$, $\rank_B M=\kappa$
and $1,\xi_1,\ldots,\xi_\kappa$ independent over $B$.
For $P\in\mathcal A_d$ of degree $m$, let $P_j$ be its homogeneous part
of degree $j$. The degree $m-1$ part of $P(Z-d)$ is
\[
 P_{m-1}(Z)-d_0\cdot\nabla P_m(Z)
 -\sum_j\xi_jM_j\cdot\nabla P_m(Z).
\]
Its coefficients lie in $B$, so each $M_j\cdot\nabla P_m=0$.
Thus $P_m$ is invariant under translation by $\operatorname{im}M$ and
$P_m(Z-d)=P_m(Z-d_0)\in B[Z]$. Subtract $P_m$ and induct on degree.
Conversely, this invariance makes $P(Z-d)=P(Z-d_0)\in B[Z]$.
Coordinates adapted to $\operatorname{im}M$ therefore give
$\mathcal A_d=B[JZ]$. Exact invariance additionally requires invariance
under $d_0$. Iteration at integer multiples, followed by polynomial
dependence on the translation parameter, gives invariance under $Bd_0$.
The joint invariant algebra is $B[J_0Z]$.

Since $Jd\in B^{N-\kappa}$, $B[JV]=B[JX]$ and its transcendence
degree is $N-\kappa$.
Since $d=A\mathbf c$ and $A$ is algebraic, $\kappa\le\rank_B A=\rank_{\Qbar}A$.
If the rows of a matrix $Q$ lie in
$\operatorname{row}(J)$, then $Qd\in B^s$ and
$\trdeg_B B(QV)=\rank_BQ$. Its relation ideal is generated by
the linear row dependencies $\mu W$, $\mu Q=0$: translation by $Qd$
preserves these equations because $\mu Qd=0$.
In particular $J_0V=J_0X$ and
$\trdeg_B B(J_0V)=\rank_BJ_0$.
The actual rank $\kappa$ and matrix $J$ are not claimed effective.
For $J_A$ whose rows span the left kernel $\{\lambda\in B^{1\times N}:\lambda A=0\}$, the uniform translation-invariant
polynomial and rational algebras are $B[J_AZ]$ and $B(J_AZ)$.
The actual row space $\{\lambda\in B^{1\times N}:\lambda d\in B\}$,
spanned by the rows of $J$, contains the left kernel of $A$ and may give
the larger coefficient-cancellation algebra $\mathcal A_d=B[JZ]$.
Once this row space is fixed, nonlinear relations among the entries of $d$
do not enlarge $\mathcal A_d$. Exact translation invariants are $B[J_0Z]$.

\subsection*{Computation for the mixed example}
In main-text Example~\ref{ex:ct2-mixed}, choose
$x=1$, $(u,m)=(2,0)$ and $(t,a)=(1,1/2)$.
Gauss's formula gives $G_{1/2}=-2\log2$; its cotangent term vanishes
and its sine sum is empty. All node logarithms and both recurrence
corrections vanish. Thus
\[
\begin{gathered}
 \mathbf c=\binom\gamma{\log2},\quad
 A=\begin{pmatrix}1&0\\1&0\\1&-2\end{pmatrix},\\
 X=\left(\Ein(1),\frac{\mathsf S(2)}{I_0(2)},
 \frac{S_{1/2}(1)}{f_{1/2}(1)}\right)^t,\quad V=X-A\mathbf c.
\end{gathered}
\]
No additive $\pi$ entry is needed, despite the multiplicative
$\Gamma(1/2)$ in the third connection value.
The first two constants cancel:
\[
 E_1(1)-\frac{K_0(2)}{I_0(2)}
 =\Ein(1)-\frac{\mathsf S(2)}{I_0(2)}=X_1-X_2.
\]
Main-text Theorem~\ref{thm:connection}, via Theorem~\ref{thm:values},
gives independence of $X_1,X_2$ over $B_*$. Thus this difference is
transcendental over $B_*$, hence over $\mathcal E$, since a polynomial
relation would contradict that independence.

\subsection*{What the actual-number bound asserts}
Main-text Theorem~\ref{thm:connection} gives
$\trdeg_F F(\mathbf c,V)\ge N$ for $F=\mathcal E,B$.
It assumes neither independence of $\mathbf c$ nor independence of $X$ over $F(\mathbf c)$.
For example, over an algebraically closed field $F$, take independent
$X_1,\ldots,X_N$, $\mathbf c=(X_1)$ and $d=(X_1,0,\ldots,0)$.
Then $V=(0,X_2,\ldots,X_N)$ and the combined degree is $N$,
although $\mathbf c$ is transcendental. Thus the argument does not give the
additive bound $N+\trdeg_F F(\mathbf c)$.
The mixed example is a simultaneous consequence of the joint value
theorem. Separate one-family statements alone do not yield independence
after all regular pairs have been adjoined. No priority claim is made.

\clearpage
\bibliographystyle{amsplain}
\label{page:references-start}
\bibliography{refs}
\label{page:references-end}
\label{page:main-end}

\end{document}
